\chapter{Chapter06 — Low Degree Junta Learning}
%%% generated by the blueprint pipeline from FABL.Chapter06.LearningAndTesting.LowDegreeJuntaLearning
\section{Overview}

Declarations of the Lean module \texttt{FABL.Chapter06.LearningAndTesting.LowDegreeJuntaLearning}.

\section{Declarations}

\begin{definition}[Low-degree junta cutoff]
\lean{FABL.lowDegreeJuntaCutoff}
\statementsource{Analysis-of-Boolean-Functions-by-Ryan-ODonnell}{
  pdf-ffbe9c7b3306-p081-b003,
  pdf-ffbe9c7b3306-p156-b002,
  pdf-ffbe9c7b3306-p156-b007,
  pdf-ffbe9c7b3306-p156-b008,
  pdf-ffbe9c7b3306-p167-b005
}
\label{FABL.lowDegreeJuntaCutoff}
\leanok
For a natural number $k$, define the cutoff to be $\lfloor 3k/4\rfloor$.
\end{definition}

\begin{definition}[Low-degree junta polynomial degree]
\lean{FABL.lowDegreeJuntaPolynomialDegree}
\statementsource{Analysis-of-Boolean-Functions-by-Ryan-ODonnell}{
  pdf-ffbe9c7b3306-p156-b007,
  pdf-ffbe9c7b3306-p167-b005
}
\label{FABL.lowDegreeJuntaPolynomialDegree}
\leanok
\uses{FABL.lowDegreeJuntaCutoff}
For each natural number $k$, define the low-degree junta polynomial degree to be
$k-c(k)$, where $c(k)$ denotes the low-degree junta cutoff at $k$.
\end{definition}

\begin{theorem}[Low-Degree Junta Cutoff Bound]
\lean{FABL.lowDegreeJuntaCutoff_le}
\label{FABL.lowDegreeJuntaCutoff_le}
\leanok
\uses{FABL.lowDegreeJuntaCutoff}
\difficulty{2}
For every natural number $k$, the cutoff $rac{3k}{4}$ rounded down satisfies
$\left\lfloor\frac{3k}{4}\right\rfloor\le k$.
\end{theorem}

\begin{theorem}[Positivity of the Polynomial Degree]
\lean{FABL.lowDegreeJuntaPolynomialDegree_pos}
\label{FABL.lowDegreeJuntaPolynomialDegree_pos}
\leanok
\uses{FABL.lowDegreeJuntaCutoff, FABL.lowDegreeJuntaPolynomialDegree}
\difficulty{2}
For every natural number $k>0$, the degree $k-\lfloor3k/4\rfloor$ is positive.
\end{theorem}

\begin{theorem}[Degree bound from the junta cutoff]
\lean{FABL.three_mul_lowDegreeJuntaPolynomialDegree_le_cutoff_add_three}
\label{FABL.three_mul_lowDegreeJuntaPolynomialDegree_le_cutoff_add_three}
\leanok
\uses{FABL.lowDegreeJuntaCutoff, FABL.lowDegreeJuntaPolynomialDegree}
\difficulty{2}
For every natural number $k$, three times the low-degree junta polynomial degree at
cutoff parameter $k$ is at most the low-degree junta cutoff at $k$ plus three.
\end{theorem}

\begin{theorem}[Cutoff and Polynomial Degree Bound]
\lean{FABL.lowDegreeJuntaCutoff_le_three_mul_polynomialDegree}
\label{FABL.lowDegreeJuntaCutoff_le_three_mul_polynomialDegree}
\leanok
\uses{FABL.lowDegreeJuntaCutoff, FABL.lowDegreeJuntaPolynomialDegree}
\difficulty{2}
For every natural number $k$, let $r=\left\lfloor\frac{3k}{4}\right\rfloor$ and $d=k-r$.
Then $r\le 3d$.
\end{theorem}

\begin{definition}[Dividing a positive learning parameter]
\lean{FABL.dividePositiveLearningParameter}
\label{FABL.dividePositiveLearningParameter}
\leanok
\uses{FABL.PositiveLearningParameter}
Given a rational failure parameter in $(0,1/2]$ and a positive natural number $d$, the
divided parameter is the original parameter divided by $d$; it again lies in $(0,1/2]$.
\end{definition}

\begin{definition}[Half failure parameter]
\lean{FABL.lowDegreeJuntaHalfFailure}
\statementsource{Analysis-of-Boolean-Functions-by-Ryan-ODonnell}{pdf-ffbe9c7b3306-p168-b001}
\label{FABL.lowDegreeJuntaHalfFailure}
\leanok
\uses{FABL.PositiveLearningParameter, FABL.dividePositiveLearningParameter}
For a positive learning parameter $\varepsilon\in(0,1/2]$, define its half-failure
parameter to be $\varepsilon/2$.
\end{definition}

\begin{definition}[Quarter failure parameter]
\lean{FABL.lowDegreeJuntaQuarterFailure}
\statementsource{Analysis-of-Boolean-Functions-by-Ryan-ODonnell}{pdf-ffbe9c7b3306-p168-b001}
\label{FABL.lowDegreeJuntaQuarterFailure}
\leanok
\uses{FABL.PositiveLearningParameter, FABL.dividePositiveLearningParameter}
For a positive learning parameter $\varepsilon\in(0,1/2]$, define its quarter-failure
parameter to be $\varepsilon/4$.
\end{definition}

\begin{theorem}[Halved Failure Parameter Value]
\lean{FABL.lowDegreeJuntaHalfFailure_value}
\label{FABL.lowDegreeJuntaHalfFailure_value}
\leanok
\uses{FABL.PositiveLearningParameter, FABL.lowDegreeJuntaHalfFailure}
\difficulty{1}
For every rational failure parameter $\delta$ with $0<\delta\le1/2$, the value of its
halved failure parameter is $\delta/2$.
\end{theorem}

\begin{theorem}[Quartered Failure Parameter]
\lean{FABL.lowDegreeJuntaQuarterFailure_value}
\label{FABL.lowDegreeJuntaQuarterFailure_value}
\leanok
\uses{FABL.PositiveLearningParameter, FABL.lowDegreeJuntaQuarterFailure}
\difficulty{1}
For every failure parameter $\delta\in\mathbb{Q}$ with $0<\delta\le1/2$, the parameter
obtained by dividing it by four has value $\delta/4$.
\end{theorem}

\begin{definition}[Junta coefficient accuracy]
\lean{FABL.lowDegreeJuntaCoefficientAccuracy}
\statementsource{Analysis-of-Boolean-Functions-by-Ryan-ODonnell}{
  pdf-ffbe9c7b3306-p080-b002,
  pdf-ffbe9c7b3306-p167-b005
}
\label{FABL.lowDegreeJuntaCoefficientAccuracy}
\leanok
\uses{FABL.PositiveLearningParameter}
For each natural number $k$, define the coefficient accuracy to be the rational number
$1/(3\cdot 2^k)$, which lies in $(0,1/2]$.
\end{definition}

\begin{definition}[Junta detection threshold]
\lean{FABL.lowDegreeJuntaDetectionThreshold}
\statementsource{Analysis-of-Boolean-Functions-by-Ryan-ODonnell}{
  pdf-ffbe9c7b3306-p083-b002,
  pdf-ffbe9c7b3306-p167-b005
}
\label{FABL.lowDegreeJuntaDetectionThreshold}
\leanok
For each natural number $k$, define the detection threshold to be the rational number
$1/(2\cdot 2^k)$.
\end{definition}

\begin{theorem}[Coefficient accuracy below the detection threshold]
\lean{FABL.lowDegreeJuntaAccuracy_lt_detectionThreshold}
\label{FABL.lowDegreeJuntaAccuracy_lt_detectionThreshold}
\leanok
\uses{FABL.lowDegreeJuntaCoefficientAccuracy, FABL.lowDegreeJuntaDetectionThreshold}
\difficulty{3}
For every natural number $k$, the coefficient accuracy $1/(3\cdot 2^k)$ is strictly less
than the detection threshold $1/(2\cdot 2^k)$.
\end{theorem}

\begin{definition}[Junta restriction dimension]
\lean{FABL.juntaRestrictionDimension}
\statementsource{Analysis-of-Boolean-Functions-by-Ryan-ODonnell}{
  pdf-ffbe9c7b3306-p075-b002,
  pdf-ffbe9c7b3306-p168-b001
}
\label{FABL.juntaRestrictionDimension}
\leanok
\uses{FABL.JuntaFreeIndex}
Let $n$ be a natural number and let $P$ be a set of coordinates in $\{0,\ldots,n-1\}$.
The junta restriction dimension is the number of coordinates outside $P$.
\end{definition}

\begin{definition}[Reindexed junta restriction]
\lean{FABL.reindexedJuntaRestriction}
\statementsource{Analysis-of-Boolean-Functions-by-Ryan-ODonnell}{pdf-ffbe9c7b3306-p168-b001}
\label{FABL.reindexedJuntaRestriction}
\leanok
\uses{FABL.BooleanFunction, FABL.JuntaFixedAssignment, FABL.fixedSignCubeEquiv, FABL.juntaRestriction, FABL.juntaRestrictionDimension}
Let $n$ be a natural number, let $f:\{-1,1\}^n\to\{-1,1\}$ be a Boolean function, and
let $P\subseteq\{0,\ldots,n-1\}$ be a coordinate set with a fixed sign assignment $z$ on
$P$. The reindexed restriction is the Boolean function on the coordinates outside $P$
obtained by fixing the coordinates in $P$ according to $z$ and identifying the remaining
coordinates with a cube of dimension equal to their number.
\end{definition}

\begin{theorem}[Evaluation of a Reindexed Junta Restriction]
\lean{FABL.reindexedJuntaRestriction_apply}
\label{FABL.reindexedJuntaRestriction_apply}
\leanok
\uses{FABL.BooleanFunction, FABL.JuntaFixedAssignment, FABL.fixedSignCubeEquiv, FABL.juntaRestriction, FABL.juntaRestrictionDimension, FABL.reindexedJuntaRestriction}
\difficulty{1}
Let $n$ be a natural number, let $f\colon\{-1,1\}^n\to\{-1,1\}$, let $P$ be a set of
coordinates, and assign a sign to each coordinate in $P$. For every sign vector $y$
indexed by the coordinates outside $P$, the reindexed restriction of $f$ at $y$ equals
the restriction evaluated on the assignment obtained by identifying those outside
coordinates with the corresponding coordinates of $y$.
\end{theorem}

\begin{definition}[Reindexing a junta restriction sample]
\lean{FABL.reindexedJuntaRestrictionExample}
\label{FABL.reindexedJuntaRestrictionExample}
\leanok
\uses{FABL.JuntaFreeAssignment, FABL.Sign, FABL.fixedSignCubeEquiv, FABL.juntaRestrictionDimension}
Let $n$ be a natural number, let $P$ be a set of coordinates in $\text{Fin}(n)$, and let
a sample consist of a sign assignment to the coordinates outside $P$ together with a
sign vector of length $n$. Reindex the first assignment as a sign vector of length equal
to the number of coordinates outside $P$, and retain the second sign vector unchanged.
\end{definition}

\begin{definition}[Low-degree junta Fourier family]
\lean{FABL.lowDegreeJuntaFourierFamily}
\statementsource{Analysis-of-Boolean-Functions-by-Ryan-ODonnell}{
  pdf-ffbe9c7b3306-p081-b003,
  pdf-ffbe9c7b3306-p156-b006,
  pdf-ffbe9c7b3306-p156-b007
}
\label{FABL.lowDegreeJuntaFourierFamily}
\leanok
\uses{FABL.juntaRestrictionDimension, FABL.lowDegreeFourierFamily, FABL.lowDegreeJuntaCutoff}
Let $n$ be a natural number, let $P$ be a set of coordinates in $\{0,\ldots,n-1\}$, and
let $k$ be a natural number. The low-degree junta Fourier family is the family of
Fourier-coordinate sets produced in dimension equal to the number of coordinates outside
$P$, using the low-degree cutoff at $k$.
\end{definition}

\begin{theorem}[Nonemptiness of the Low-Degree Fourier Family]
\lean{FABL.lowDegreeJuntaFourierFamily_nonempty}
\label{FABL.lowDegreeJuntaFourierFamily_nonempty}
\leanok
\uses{FABL.lowDegreeFourierFamily_nonempty, FABL.lowDegreeJuntaFourierFamily}
\difficulty{1}
Let $n,k$ be natural numbers and let $P$ be a subset of an $n$-element index set. The
family of subsets of the indices outside $P$ having size at most $\lfloor3k/4\rfloor$ is
nonempty.
\end{theorem}

\begin{theorem}[Membership in the low-degree junta Fourier family]
\lean{FABL.mem_lowDegreeJuntaFourierFamily}
\label{FABL.mem_lowDegreeJuntaFourierFamily}
\leanok
\uses{FABL.juntaRestrictionDimension, FABL.lowDegreeJuntaCutoff, FABL.lowDegreeJuntaFourierFamily, FABL.mem_lowDegreeFourierFamily}
\difficulty{1}
Let $n$ and $k$ be natural numbers, let $P$ be a set of coordinates, and let $S$ be a
subset of the coordinates outside $P$, represented in the corresponding restriction
dimension. Then $S$ belongs to the low-degree Fourier family for that restriction if and
only if $|S|\\le\\lfloor 3k/4\\rfloor$.
\end{theorem}

\begin{definition}[Coefficient confidence parameter]
\lean{FABL.lowDegreeJuntaCoefficientConfidence}
\statementsource{Analysis-of-Boolean-Functions-by-Ryan-ODonnell}{
  pdf-ffbe9c7b3306-p080-b002,
  pdf-ffbe9c7b3306-p156-b006,
  pdf-ffbe9c7b3306-p167-b005
}
\label{FABL.lowDegreeJuntaCoefficientConfidence}
\leanok
\uses{FABL.PositiveLearningParameter, FABL.dividePositiveLearningParameter, FABL.lowDegreeJuntaFourierFamily, FABL.lowDegreeJuntaFourierFamily_nonempty, FABL.lowDegreeJuntaQuarterFailure}
Let $n$ and $k$ be nonnegative integers, let $P$ be a set of coordinates in $\text{Fin
}n$, and let $ailure\text{-parameter}$ be a rational number in $(0,1/2]$. Form the
family of subsets of the coordinates outside $P$ having size at most $rac{3k}{4}$, with
the cutoff rounded down, and let $N$ be its cardinality. The coefficient confidence
parameter is the failure parameter divided by $4N$.
\end{definition}

\begin{theorem}[Coefficient confidence value]
\lean{FABL.lowDegreeJuntaCoefficientConfidence_value}
\label{FABL.lowDegreeJuntaCoefficientConfidence_value}
\leanok
\uses{FABL.PositiveLearningParameter, FABL.lowDegreeJuntaCoefficientConfidence, FABL.lowDegreeJuntaFourierFamily}
\difficulty{1}
Let $n$ be a natural number, let $P\subseteq[n]$, let $k$ be a natural number, and let
$\delta\in\mathbb{Q}$ satisfy $0<\delta\leq 1/2$. Write $\mathcal F$ for the family of
subsets of the coordinates outside $P$ having size at most $\lfloor 3k/4\rfloor$. The
coefficient confidence parameter has value $\delta/(4|\mathcal F|)$.
\end{theorem}

\begin{definition}[Coefficient estimation sample count]
\lean{FABL.lowDegreeJuntaCoefficientSampleCount}
\statementsource{Analysis-of-Boolean-Functions-by-Ryan-ODonnell}{
  pdf-ffbe9c7b3306-p080-b003,
  pdf-ffbe9c7b3306-p156-b006,
  pdf-ffbe9c7b3306-p167-b005
}
\label{FABL.lowDegreeJuntaCoefficientSampleCount}
\leanok
\uses{FABL.PositiveLearningParameter, FABL.fourierEstimatorSampleCount, FABL.lowDegreeJuntaCoefficientAccuracy, FABL.lowDegreeJuntaCoefficientConfidence}
Let $n$ and $k$ be nonnegative integers, let $P$ be a set of coordinates in $\text{Fin
}n$, and let the failure parameter be rational and lie in $(0,1/2]$. Set the accuracy to
$1/(3\cdot 2^k)$ and set the confidence parameter to the coefficient confidence
parameter for $P$, $k$, and the given failure parameter. The sample count is the ceiling
of $2b/\varepsilon^2$, where $b$ is the least integer $m$ such that $2^m\geq \lceil
2/\delta\rceil$, with $\varepsilon$ the accuracy and $\delta$ the confidence parameter.
\end{definition}

\begin{definition}[Junta polynomial sample count]
\lean{FABL.lowDegreeJuntaPolynomialSampleCount}
\statementsource{Analysis-of-Boolean-Functions-by-Ryan-ODonnell}{
  pdf-ffbe9c7b3306-p081-b004,
  pdf-ffbe9c7b3306-p156-b007,
  pdf-ffbe9c7b3306-p167-b005
}
\label{FABL.lowDegreeJuntaPolynomialSampleCount}
\leanok
\uses{FABL.PositiveLearningParameter, FABL.juntaRestrictionDimension, FABL.lowDegreeF₂LearningSampleCount, FABL.lowDegreeJuntaPolynomialDegree, FABL.lowDegreeJuntaQuarterFailure}
Let $n$ and $k$ be nonnegative integers, let $P$ be a set of coordinates in $\text{Fin
}n$, and let the failure parameter be rational and lie in $(0,1/2]$. Let $d$ be the
number of coordinates outside $P$, let $r=k-\lfloor 3k/4\rfloor$, and let $\delta$ be
one quarter of the failure parameter. The sample count is $2^r$ times the sum of the
number of subsets of a $d$-element set of size at most $r$ and the least integer $m$
such that $2^m\geq\lceil 2/\delta\rceil$.
\end{definition}

\begin{definition}[Restricted junta sample count]
\lean{FABL.lowDegreeJuntaRestrictedSampleCount}
\statementsource{Analysis-of-Boolean-Functions-by-Ryan-ODonnell}{pdf-ffbe9c7b3306-p168-b001}
\label{FABL.lowDegreeJuntaRestrictedSampleCount}
\leanok
\uses{FABL.PositiveLearningParameter, FABL.lowDegreeJuntaCoefficientSampleCount, FABL.lowDegreeJuntaFourierFamily, FABL.lowDegreeJuntaPolynomialSampleCount}
Let $n$ and $k$ be nonnegative integers, let $P$ be a set of coordinates in $\text{Fin
}n$, and let the failure parameter be rational and lie in $(0,1/2]$. The restricted
sample count is the cardinality of the family of subsets of the coordinates outside $P$
of size at most $\lfloor 3k/4\rfloor$, multiplied by the coefficient estimation sample
count, plus the junta polynomial sample count.
\end{definition}

\begin{definition}[Rank of a low-degree junta frequency]
\lean{FABL.lowDegreeJuntaFrequencyRank}
\label{FABL.lowDegreeJuntaFrequencyRank}
\leanok
\uses{FABL.lowDegreeJuntaFourierFamily}
Let $n$ and $k$ be natural numbers, let $P$ be a set of coordinates in $\text{Fin}(n)$,
and let $S$ be a member of the family of subsets of the coordinates outside $P$ having
size at most $rac{3k}{4}$ rounded down. Its rank is its position in the canonical
enumeration of that finite family, represented by an element of $\text{Fin}$ whose size
is the cardinality of the family.
\end{definition}

\begin{definition}[Index of a coefficient-estimation sample]
\lean{FABL.lowDegreeJuntaCoefficientSampleIndex}
\label{FABL.lowDegreeJuntaCoefficientSampleIndex}
\leanok
\uses{FABL.PositiveLearningParameter, FABL.lowDegreeJuntaCoefficientSampleCount, FABL.lowDegreeJuntaFourierFamily, FABL.lowDegreeJuntaFrequencyRank, FABL.lowDegreeJuntaPolynomialSampleCount, FABL.lowDegreeJuntaRestrictedSampleCount}
Let $n,k$ be natural numbers, let $P$ be a set of coordinates in $\text{Fin}(n)$, and let
the failure parameter lie in $(0,1/2]$. For a frequency $S$ in the low-degree junta
Fourier family and a coefficient-estimation sample index $i$, assign the index
$\text{rank}(S)m+i$ in the combined restricted sample space, where $m$ is the
coefficient-estimation sample count.
\end{definition}

\begin{definition}[Index of a polynomial-estimation sample]
\lean{FABL.lowDegreeJuntaPolynomialSampleIndex}
\label{FABL.lowDegreeJuntaPolynomialSampleIndex}
\leanok
\uses{FABL.PositiveLearningParameter, FABL.lowDegreeJuntaCoefficientSampleCount, FABL.lowDegreeJuntaFourierFamily, FABL.lowDegreeJuntaPolynomialSampleCount, FABL.lowDegreeJuntaRestrictedSampleCount}
Let $n,k$ be natural numbers, let $P$ be a set of coordinates in $\text{Fin}(n)$, and let
the failure parameter lie in $(0,1/2]$. A polynomial-estimation sample with index $i$ is
assigned the index $N m+i$ in the combined restricted sample space, where $N$ is the
size of the low-degree junta Fourier family and $m$ is the coefficient-estimation sample
count.
\end{definition}

\begin{definition}[Embedding coefficient-estimation sample indices]
\lean{FABL.lowDegreeJuntaCoefficientSampleEmbedding}
\label{FABL.lowDegreeJuntaCoefficientSampleEmbedding}
\leanok
\uses{FABL.PositiveLearningParameter, FABL.lowDegreeJuntaCoefficientSampleCount, FABL.lowDegreeJuntaCoefficientSampleIndex, FABL.lowDegreeJuntaFourierFamily, FABL.lowDegreeJuntaRestrictedSampleCount}
Let $n,k$ be natural numbers, let $P$ be a set of coordinates in $\text{Fin}(n)$, and let
the failure parameter lie in $(0,1/2]$. For each frequency $S$ in the low-degree junta
Fourier family, the coefficient-estimation sample indices embed injectively into the
combined restricted sample indices by sending each index to its block, indexed by the
rank of $S$.
\end{definition}

\begin{definition}[Embedding polynomial-estimation sample indices]
\lean{FABL.lowDegreeJuntaPolynomialSampleEmbedding}
\label{FABL.lowDegreeJuntaPolynomialSampleEmbedding}
\leanok
\uses{FABL.PositiveLearningParameter, FABL.lowDegreeJuntaPolynomialSampleCount, FABL.lowDegreeJuntaPolynomialSampleIndex, FABL.lowDegreeJuntaRestrictedSampleCount}
Let $n,k$ be natural numbers, let $P$ be a set of coordinates in $\text{Fin}(n)$, and let
the failure parameter lie in $(0,1/2]$. Polynomial-estimation sample indices embed
injectively into the combined restricted sample indices by adding the total number of
coefficient-estimation samples across the low-degree junta Fourier family.
\end{definition}

\begin{definition}[Batch equivalence for free inputs]
\lean{FABL.juntaFreeInputBatchEquiv}
\label{FABL.juntaFreeInputBatchEquiv}
\leanok
\uses{FABL.JuntaFreeAssignment, FABL.fixedSignCubeEquiv, FABL.juntaRestrictionDimension}
Let $n$ be a natural number, let $P$ be a set of coordinates in $\\{0,\\ldots,n-1\\}$,
and let $M$ be a natural number. There is a bijection between batches of $M$ assignments
to the coordinates outside $P$ and batches of $M$ sign vectors of length equal to the
number of coordinates outside $P$, obtained by converting each assignment to its
sign-vector representation coordinatewise.
\end{definition}

\begin{definition}[Low-degree junta coefficient input row]
\lean{FABL.lowDegreeJuntaCoefficientInputRow}
\label{FABL.lowDegreeJuntaCoefficientInputRow}
\leanok
\uses{FABL.JuntaFreeAssignment, FABL.PositiveLearningParameter, FABL.fixedSignCubeEquiv, FABL.juntaRestrictionDimension, FABL.lowDegreeJuntaCoefficientSampleCount, FABL.lowDegreeJuntaCoefficientSampleEmbedding, FABL.lowDegreeJuntaFourierFamily, FABL.lowDegreeJuntaRestrictedSampleCount}
Let $n$ and $k$ be natural numbers, let $P$ be a subset of the $n$ coordinates, and let
the failure parameter lie in $(0,1/2]$. For a frequency $S$ in the low-degree Fourier
family on the coordinates outside $P$ and a sequence of restricted assignments of the
prescribed restricted sample count, the coefficient input row is the sequence, of length
equal to the coefficient sample count, of sign vectors obtained by converting each
assignment at its designated sample position for $S$ to the sign cube on the restricted
coordinates.
\end{definition}

\begin{definition}[Polynomial estimation input row]
\lean{FABL.lowDegreeJuntaPolynomialInputRow}
\label{FABL.lowDegreeJuntaPolynomialInputRow}
\leanok
\uses{FABL.JuntaFreeAssignment, FABL.PositiveLearningParameter, FABL.fixedSignCubeEquiv, FABL.juntaRestrictionDimension, FABL.lowDegreeJuntaPolynomialSampleCount, FABL.lowDegreeJuntaPolynomialSampleEmbedding, FABL.lowDegreeJuntaRestrictedSampleCount}
Let $n$ be a natural number, let $P$ be a set of coordinates in $0,\\ldots,n-1$, let
$k$ be a natural number, and let the failure parameter lie in $(0,1/2]$. Given a batch
of the prescribed restricted sample size, the polynomial estimation input row is the
sequence of sign vectors obtained by converting the inputs at the sample positions
assigned to polynomial estimation.
\end{definition}

\begin{theorem}[Uniform distribution of junta coefficient inputs]
\lean{FABL.map_uniformPMF_lowDegreeJuntaCoefficientInputRow}
\label{FABL.map_uniformPMF_lowDegreeJuntaCoefficientInputRow}
\leanok
\uses{FABL.JuntaFreeAssignment, FABL.PositiveLearningParameter, FABL.juntaFreeInputBatchEquiv, FABL.juntaRestrictionDimension, FABL.lowDegreeJuntaCoefficientInputRow, FABL.lowDegreeJuntaCoefficientSampleCount, FABL.lowDegreeJuntaCoefficientSampleEmbedding, FABL.lowDegreeJuntaFourierFamily, FABL.lowDegreeJuntaRestrictedSampleCount, FABL.map_uniformPMF_equiv, FABL.map_uniformPMF_injection_projection, FABL.uniformPMF}
\difficulty{4}
Let $n,k$ be natural numbers, let $P\subseteq\{0,\ldots,n-1\}$, let $\delta$ be a
rational number with $0<\delta\leq 1/2$, and choose a member $S$ of the family of
subsets of the free coordinates of size at most the cutoff associated with $k$. If a
batch of restricted assignments is chosen uniformly, then extracting the
coefficient-input rows indexed by $S$ gives a uniformly distributed batch of sign
vectors of the prescribed coefficient sample count, with one sign for each free
coordinate.
\end{theorem}

\begin{theorem}[Uniform distribution of junta polynomial inputs]
\lean{FABL.map_uniformPMF_lowDegreeJuntaPolynomialInputRow}
\label{FABL.map_uniformPMF_lowDegreeJuntaPolynomialInputRow}
\leanok
\uses{FABL.JuntaFreeAssignment, FABL.PositiveLearningParameter, FABL.juntaFreeInputBatchEquiv, FABL.juntaRestrictionDimension, FABL.lowDegreeJuntaPolynomialInputRow, FABL.lowDegreeJuntaPolynomialSampleCount, FABL.lowDegreeJuntaPolynomialSampleEmbedding, FABL.lowDegreeJuntaRestrictedSampleCount, FABL.map_uniformPMF_equiv, FABL.map_uniformPMF_injection_projection, FABL.uniformPMF}
\difficulty{4}
Let $n,k$ be natural numbers, let $P\subseteq\{0,\ldots,n-1\}$, and let $\delta$ be a
rational number with $0<\delta\leq 1/2$. If a batch of restricted assignments is chosen
uniformly, then extracting the polynomial-input rows gives a uniformly distributed batch
of sign vectors, with one sign for each free coordinate and with batch size equal to the
polynomial sample count determined by $P$, $k$, and $\delta$.
\end{theorem}

\begin{definition}[Junta coefficient sample sequence]
\lean{FABL.lowDegreeJuntaCoefficientSamples}
\statementsource{Analysis-of-Boolean-Functions-by-Ryan-ODonnell}{
  pdf-ffbe9c7b3306-p080-b003,
  pdf-ffbe9c7b3306-p156-b006,
  pdf-ffbe9c7b3306-p167-b005
}
\label{FABL.lowDegreeJuntaCoefficientSamples}
\leanok
\uses{FABL.JuntaFreeAssignment, FABL.PositiveLearningParameter, FABL.Sign, FABL.juntaRestrictionDimension, FABL.lowDegreeJuntaCoefficientSampleCount, FABL.lowDegreeJuntaCoefficientSampleIndex, FABL.lowDegreeJuntaFourierFamily, FABL.lowDegreeJuntaRestrictedSampleCount, FABL.reindexedJuntaRestrictionExample}
Let $n$ and $k$ be nonnegative integers, let $P$ be a set of coordinates in $\text{Fin
}n$, and let the failure parameter be rational and lie in $(0,1/2]$. Given a batch of
restricted examples, indexed by the restricted sample count and consisting of an
assignment of signs to the coordinates outside $P$ together with a sign label, and a
subset $S$ in the low-degree Fourier family, the coefficient sample sequence has length
equal to the coefficient estimation sample count. Its entries are the reindexed
restriction examples selected from the batch by the prescribed coefficient sample
indices for $S$.
\end{definition}

\begin{definition}[Junta polynomial sample sequence]
\lean{FABL.lowDegreeJuntaPolynomialSamples}
\statementsource{Analysis-of-Boolean-Functions-by-Ryan-ODonnell}{
  pdf-ffbe9c7b3306-p081-b004,
  pdf-ffbe9c7b3306-p156-b007,
  pdf-ffbe9c7b3306-p167-b005
}
\label{FABL.lowDegreeJuntaPolynomialSamples}
\leanok
\uses{FABL.JuntaFreeAssignment, FABL.PositiveLearningParameter, FABL.Sign, FABL.juntaRestrictionDimension, FABL.lowDegreeJuntaPolynomialSampleCount, FABL.lowDegreeJuntaPolynomialSampleIndex, FABL.lowDegreeJuntaRestrictedSampleCount, FABL.reindexedJuntaRestrictionExample}
Let $n$ and $k$ be nonnegative integers, let $P$ be a set of coordinates in $\text{Fin
}n$, and let the failure parameter be rational and lie in $(0,1/2]$. Given a batch of
restricted examples indexed by the restricted sample count, the polynomial sample
sequence has length equal to the junta polynomial sample count. Its entries are the
reindexed restriction examples selected from the batch by the prescribed polynomial
sample indices.
\end{definition}

\begin{definition}[Junta Fourier coefficient estimate]
\lean{FABL.lowDegreeJuntaCoefficientEstimate}
\statementsource{Analysis-of-Boolean-Functions-by-Ryan-ODonnell}{
  pdf-ffbe9c7b3306-p080-b002,
  pdf-ffbe9c7b3306-p080-b003,
  pdf-ffbe9c7b3306-p156-b006
}
\label{FABL.lowDegreeJuntaCoefficientEstimate}
\leanok
\uses{FABL.JuntaFreeAssignment, FABL.PositiveLearningParameter, FABL.Sign, FABL.empiricalFourierCoeff, FABL.lowDegreeJuntaCoefficientSamples, FABL.lowDegreeJuntaFourierFamily, FABL.lowDegreeJuntaRestrictedSampleCount}
Let $n$ and $k$ be nonnegative integers, let $P$ be a set of coordinates in $\text{Fin
}n$, and let the failure parameter be rational and lie in $(0,1/2]$. Given a batch of
restricted examples and a member $S$ of the low-degree junta Fourier family, the
coefficient estimate is the empirical Fourier coefficient at $S$ computed from the
coefficient sample sequence for that batch and $S$; it is rational-valued.
\end{definition}

\begin{definition}[Detected nonempty Fourier frequencies]
\lean{FABL.lowDegreeJuntaDetectedFrequencies}
\statementsource{Analysis-of-Boolean-Functions-by-Ryan-ODonnell}{
  pdf-ffbe9c7b3306-p083-b002,
  pdf-ffbe9c7b3306-p156-b006
}
\label{FABL.lowDegreeJuntaDetectedFrequencies}
\leanok
\uses{FABL.JuntaFreeAssignment, FABL.PositiveLearningParameter, FABL.Sign, FABL.juntaRestrictionDimension, FABL.lowDegreeJuntaCoefficientEstimate, FABL.lowDegreeJuntaDetectionThreshold, FABL.lowDegreeJuntaFourierFamily, FABL.lowDegreeJuntaRestrictedSampleCount}
Let $n$ and $k$ be nonnegative integers, let $P$ be a set of coordinates in $\text{Fin
}n$, and let the failure parameter be rational and lie in $(0,1/2]$. Given a batch of
restricted examples, the detected frequencies are those members $S$ of the low-degree
junta Fourier family that are nonempty and whose coefficient estimate, computed using
$S$ as a member of that family, has absolute value at least the detection threshold for
$k$.
\end{definition}

\begin{definition}[First detected low-degree junta frequency]
\lean{FABL.firstLowDegreeJuntaDetectedFrequency}
\label{FABL.firstLowDegreeJuntaDetectedFrequency}
\leanok
\uses{FABL.JuntaFreeAssignment, FABL.PositiveLearningParameter, FABL.Sign, FABL.juntaRestrictionDimension, FABL.lowDegreeJuntaDetectedFrequencies, FABL.lowDegreeJuntaRestrictedSampleCount}
Let $n$ and $k$ be nonnegative integers, let $P$ be a subset of the $n$ coordinates, and
let the failure parameter $ailure$ be a rational number in $(0,1/2]$. Given a batch of
the prescribed size, whose entries pair an assignment to the coordinates outside $P$
with a sign, the construction returns no frequency when the set of detected frequencies
is empty; otherwise it returns a nonempty frequency from that set.
\end{definition}

\begin{definition}[Low degree junta polynomial hypothesis]
\lean{FABL.lowDegreeJuntaPolynomialHypothesis}
\statementsource{Analysis-of-Boolean-Functions-by-Ryan-ODonnell}{
  pdf-ffbe9c7b3306-p080-b006,
  pdf-ffbe9c7b3306-p081-b003,
  pdf-ffbe9c7b3306-p156-b007,
  pdf-ffbe9c7b3306-p167-b005
}
\label{FABL.lowDegreeJuntaPolynomialHypothesis}
\leanok
\uses{FABL.JuntaFreeAssignment, FABL.LowDegreeF₂Hypothesis, FABL.PositiveLearningParameter, FABL.Sign, FABL.juntaRestrictionDimension, FABL.lowDegreeF₂PolynomialLearnerLabeledOutput, FABL.lowDegreeJuntaPolynomialDegree, FABL.lowDegreeJuntaPolynomialSampleCount, FABL.lowDegreeJuntaPolynomialSamples, FABL.lowDegreeJuntaRestrictedSampleCount}
Let $n,k$ be nonnegative integers, let $P\subseteq\{0,\ldots,n-1\}$, and let the failure
parameter lie in $(0,1/2]$. Given a batch of the prescribed restricted sample size,
consisting of a free-coordinate sign assignment and a label for each sample, apply the
low-degree polynomial learner to the prescribed polynomial-learning portion of the
batch. The result is an $\mathbb{F}_2$-coefficient hypothesis on the free-coordinate
space, with degree bound $k-\lfloor 3k/4\rfloor$.
\end{definition}

\begin{theorem}[Ideal coefficient samples under reindexing]
\lean{FABL.lowDegreeJuntaCoefficientSamples_ideal}
\label{FABL.lowDegreeJuntaCoefficientSamples_ideal}
\leanok
\uses{FABL.BooleanFunction, FABL.JuntaFixedAssignment, FABL.JuntaFreeAssignment, FABL.PositiveLearningParameter, FABL.juntaRestrictionLabeledBatch, FABL.lowDegreeJuntaCoefficientInputRow, FABL.lowDegreeJuntaCoefficientSampleEmbedding, FABL.lowDegreeJuntaCoefficientSamples, FABL.lowDegreeJuntaFourierFamily, FABL.lowDegreeJuntaRestrictedSampleCount, FABL.reindexedJuntaRestriction, FABL.reindexedJuntaRestrictionExample}
\difficulty{1}
Let $n$ be a natural number, let $f:\{-1,1\}^{[n]}\to\{-1,1\}$ be a Boolean function,
let $P\subseteq[n]$, and fix a sign assignment $z:P\to\{-1,1\}$. Let $k$ be a natural
number, let $\delta\in\mathbb{Q}$ satisfy $0<\delta\leq1/2$, and let the input batch
consist of free-coordinate sign assignments indexed by the restricted sample count for
$(P,k,\delta)$. For every frequency $S$ among the subsets of free coordinates of size at
most $\lfloor3k/4\rfloor$, the batch of reindexed labeled samples obtained by
restricting $f$ at $z$ agrees pointwise with the pairs consisting of each corresponding
reindexed input and the value of the reindexed restricted function at that input.
\end{theorem}

\begin{theorem}[Ideal Junta Polynomial Samples]
\lean{FABL.lowDegreeJuntaPolynomialSamples_ideal}
\label{FABL.lowDegreeJuntaPolynomialSamples_ideal}
\leanok
\uses{FABL.BooleanFunction, FABL.JuntaFixedAssignment, FABL.JuntaFreeAssignment, FABL.PositiveLearningParameter, FABL.juntaRestrictionLabeledBatch, FABL.lowDegreeJuntaPolynomialInputRow, FABL.lowDegreeJuntaPolynomialSampleEmbedding, FABL.lowDegreeJuntaPolynomialSamples, FABL.lowDegreeJuntaRestrictedSampleCount, FABL.reindexedJuntaRestriction, FABL.reindexedJuntaRestrictionExample}
\difficulty{1}
Let $n,k$ be nonnegative integers, let $f:\{\pm1\}^n\to\{\pm1\}$ be a Boolean function,
let $P$ be a subset of the $n$ coordinates, and fix signs on $P$. Let $0<\delta\le1/2$,
and let a batch of restricted inputs and sign labels be indexed by the restricted sample
count. The polynomial sample list selected from this batch is exactly the list whose
$i$th entry consists of the reindexed input selected at position $i$ and the value of
the reindexed restriction of $f$ at that input; the list has the prescribed polynomial
sample count.
\end{theorem}

\begin{theorem}[Ideal coefficient estimate under reindexing]
\lean{FABL.lowDegreeJuntaCoefficientEstimate_cast_ideal}
\label{FABL.lowDegreeJuntaCoefficientEstimate_cast_ideal}
\leanok
\uses{FABL.BooleanFunction, FABL.JuntaFixedAssignment, FABL.JuntaFreeAssignment, FABL.PositiveLearningParameter, FABL.empiricalFourierCoeff_cast, FABL.juntaRestrictionLabeledBatch, FABL.lowDegreeJuntaCoefficientEstimate, FABL.lowDegreeJuntaCoefficientInputRow, FABL.lowDegreeJuntaCoefficientSamples_ideal, FABL.lowDegreeJuntaFourierFamily, FABL.lowDegreeJuntaRestrictedSampleCount, FABL.realEmpiricalFourierCoeff, FABL.reindexedJuntaRestriction}
\difficulty{3}
Let $n$ be a natural number, let $f:\{-1,1\}^{[n]}\to\{-1,1\}$ be a Boolean function,
let $P\subseteq[n]$, and fix a sign assignment $z:P\to\{-1,1\}$. Let $k$ be a natural
number, let $\delta\in\mathbb{Q}$ satisfy $0<\delta\leq1/2$, and take a batch of
free-coordinate sign assignments indexed by the restricted sample count for
$(P,k,\delta)$. For any frequency $S$ in the family of subsets of the free coordinates
of size at most $\lfloor3k/4\rfloor$, the real value of the empirical rational Fourier
coefficient computed from the batch labeled by the restriction of $f$ at $z$ equals the
empirical real Fourier coefficient of the reindexed restricted function, computed on the
corresponding reindexed inputs.
\end{theorem}

\begin{theorem}[Ideal Low-Degree Junta Hypothesis]
\lean{FABL.lowDegreeJuntaPolynomialHypothesis_ideal}
\label{FABL.lowDegreeJuntaPolynomialHypothesis_ideal}
\leanok
\uses{FABL.BooleanFunction, FABL.JuntaFixedAssignment, FABL.JuntaFreeAssignment, FABL.PositiveLearningParameter, FABL.juntaRestrictionDimension, FABL.juntaRestrictionLabeledBatch, FABL.lowDegreeF₂PolynomialLearnerLabeledOutput, FABL.lowDegreeJuntaPolynomialDegree, FABL.lowDegreeJuntaPolynomialHypothesis, FABL.lowDegreeJuntaPolynomialInputRow, FABL.lowDegreeJuntaPolynomialSampleCount, FABL.lowDegreeJuntaPolynomialSamples_ideal, FABL.lowDegreeJuntaRestrictedSampleCount, FABL.reindexedJuntaRestriction}
\difficulty{3}
Let $n,k$ be nonnegative integers, let $f:\{\pm1\}^n\to\{\pm1\}$ be a Boolean function,
let $P\subseteq\{0,\ldots,n-1\}$, and fix an assignment of signs to the coordinates in
$P$. Let $\delta\in\mathbb{Q}$ satisfy $0<\delta\le 1/2$, and choose an input batch
indexed by the restricted sample count, with each entry an assignment of signs to the
coordinates outside $P$. The low-degree $\mathbb{F}_2$ learner applied to the batch
labeled by the restriction of $f$ equals the learner applied to the corresponding
reindexed inputs labeled by that same restricted function, where the learner's degree
bound is $k-\lfloor 3k/4\rfloor$ and its sample count is the prescribed polynomial
sample count.
\end{theorem}

\begin{definition}[Boolean function associated to an $\mathbb{F}_2$ hypothesis]
\lean{FABL.f₂HypothesisBooleanFunction}
\leanok
\uses{FABL.BooleanFunction, FABL.F₂BooleanFunction, FABL.binaryCubeSignEquiv, FABL.binarySignEquiv}
\statementsource{Analysis-of-Boolean-Functions-by-Ryan-ODonnell}{
  pdf-ffbe9c7b3306-p156-b007,
  pdf-ffbe9c7b3306-p167-b005,
  pdf-ffbe9c7b3306-p169-b008
}
For a nonnegative integer $r$, let $h$ be a function from the $r$-dimensional cube over
$\mathbb{F}_2$ to $\mathbb{F}_2$. The associated Boolean function on the sign cube is
obtained by converting each sign input to its $\mathbb{F}_2$ cube point, evaluating $h$,
and converting the output to a sign.
\end{definition}

\begin{theorem}[Encoding and decoding Boolean functions]
\lean{FABL.f₂HypothesisBooleanFunction_encoding}
\leanok
\uses{FABL.BooleanFunction, FABL.booleanFunctionF₂Encoding, FABL.f₂HypothesisBooleanFunction}
\difficulty{1}
Let $r$ be a natural number and let $g$ be a Boolean function on the sign cube
$\{\pm1\}^r$ with values in $\{\pm1\}$. Converting $g$ to its binary-valued function on
$\mathbb F_2^r$ and then converting back recovers $g$ at every input.
\end{theorem}

\begin{definition}[Nonconstant nonzero low degree monomials]
\lean{FABL.nonconstantLowDegreeMonomials}
\statementsource{Analysis-of-Boolean-Functions-by-Ryan-ODonnell}{
  pdf-ffbe9c7b3306-p156-b007,
  pdf-ffbe9c7b3306-p169-b008
}
\label{FABL.nonconstantLowDegreeMonomials}
\leanok
\uses{FABL.LowDegreeF₂Coefficients, FABL.LowDegreeMonomial}
Let $r,\ell$ be nonnegative integers, and assign an $\mathbb{F}_2$ coefficient to each
subset of $\{0,\ldots,r-1\}$ of size at most $\ell$. The selected monomials are exactly
those indexed by a nonempty subset whose coefficient is nonzero.
\end{definition}

\begin{definition}[First nonconstant monomial with nonzero coefficient]
\lean{FABL.firstNonconstantLowDegreeMonomial}
\label{FABL.firstNonconstantLowDegreeMonomial}
\leanok
\uses{FABL.LowDegreeF₂Coefficients, FABL.LowDegreeMonomial, FABL.nonconstantLowDegreeMonomials}
Let $r$ and $\ell$ be natural numbers, and assign a coefficient in $\mathbb{F}_2$ to
each monomial in $r$ variables of degree at most $\ell$. The construction returns the
first, in a finite enumeration of the eligible monomials, that is nonconstant and has
nonzero coefficient; it returns none if there is no such monomial.
\end{definition}

\begin{theorem}[Invariant coordinate forces a zero coefficient]
\lean{FABL.anfCoeff_eq_zero_of_update_eq}
\label{FABL.anfCoeff_eq_zero_of_update_eq}
\leanok
\uses{FABL.F₂BooleanFunction, FABL.F₂Cube, FABL.anfCoeff, FABL.f₂CubeOfFinset, FABL.f₂CubeOfFinset_apply}
\difficulty{5}
Let $r$ be a natural number, let $f:\{0,1\}^r\to\mathbb F_2$, let
$S\subseteq\{0,\ldots,r-1\}$, and let $i\in S$. Suppose that for every input $x$, the
value of $f$ is unchanged when coordinate $i$ is set to $1$ rather than to $0$, with all
other coordinates fixed. Then the algebraic normal form coefficient of $f$ indexed by
$S$ is $0$.
\end{theorem}

\begin{theorem}[Is relevant f₂hypothesis boolean function of anf coeff ne zero]
\lean{FABL.isRelevant_f₂HypothesisBooleanFunction_of_anfCoeff_ne_zero}
\leanok
\uses{FABL.BooleanFunction.toReal, FABL.F₂BooleanFunction, FABL.F₂Cube, FABL.IsRelevant, FABL.anfCoeff, FABL.anfCoeff_eq_zero_of_update_eq, FABL.binaryCubeSignEquiv, FABL.binaryCubeSignEquiv_apply, FABL.binarySignEquiv, FABL.f₂HypothesisBooleanFunction, FABL.isRelevant_iff_exists_setCoordinate_ne, FABL.setCoordinate, FABL.setCoordinate_apply_of_ne, FABL.setCoordinate_apply_self, FABL.signEncode_one, FABL.signEncode_zero, FABL.signValue_injective}
\statementsource{Analysis-of-Boolean-Functions-by-Ryan-ODonnell}{pdf-ffbe9c7b3306-p156-b006}
\difficulty{5}
% dataset-pipeline: stub -- prose pending build_informal_statements.py
Statement pending regeneration.
\end{theorem}

\begin{theorem}[Constant evaluation with no nonconstant support]
\lean{FABL.lowDegreeF₂Eval_eq_const_of_nonconstantSupport_empty}
\leanok
\uses{FABL.LowDegreeF₂Coefficients, FABL.LowDegreeMonomial, FABL.lowDegreeF₂Eval, FABL.nonconstantLowDegreeMonomials}
\difficulty{4}
Let $r$ and $\ell$ be nonnegative integers, and let a coefficient family assign an
element of $\mathbb{F}_2$ to each subset of $\{1,\ldots,r\}$ of size at most $\ell$. If
no nonempty subset has a nonzero coefficient, then the resulting polynomial evaluation
is the constant function whose value is the coefficient of the empty set.
\end{theorem}

\begin{theorem}[Soundness of the first nonconstant monomial]
\lean{FABL.firstNonconstantLowDegreeMonomial_sound}
\label{FABL.firstNonconstantLowDegreeMonomial_sound}
\leanok
\uses{FABL.BooleanFunction.toReal, FABL.IsRelevant, FABL.LowDegreeF₂Coefficients, FABL.LowDegreeMonomial, FABL.anfCoeff_lowDegreeF₂Eval, FABL.extendLowDegreeF₂Coefficients, FABL.firstNonconstantLowDegreeMonomial, FABL.f₂HypothesisBooleanFunction, FABL.isRelevant_f₂HypothesisBooleanFunction_of_anfCoeff_ne_zero, FABL.lowDegreeF₂Eval, FABL.lowDegreeF₂Eval_eq_const_of_nonconstantSupport_empty, FABL.nonconstantLowDegreeMonomials}
\difficulty{4}
Let $r,\ell\in\mathbb{N}$, and assign a coefficient in $\mathbb{F}_2$ to each monomial
indexed by a subset of $\{0,\ldots,r-1\}$ of size at most $\ell$. If there is no
nonconstant monomial with nonzero coefficient, then the polynomial represented by these
coefficients is the constant function whose value is the coefficient of the empty
monomial. If there is such a monomial, the Boolean function obtained by evaluating the
polynomial and encoding its values as signs has positive influence in the least
coordinate of the first nonconstant monomial selected by the specified rule.
\end{theorem}

\begin{definition}[Enumeration of free coordinates]
\lean{FABL.juntaFreeCoordinateOfFin}
\label{FABL.juntaFreeCoordinateOfFin}
\leanok
\uses{FABL.FixedIndex, FABL.JuntaFreeIndex, FABL.juntaRestrictionDimension}
Let $n$ be a natural number and let $P$ be a set of coordinates in $0,\\ldots,n-1$.
The $i$-th coordinate in the enumeration of the coordinates outside $P$ is the
coordinate represented by index $i$, for $i$ ranging from zero to one less than the
number of coordinates outside $P$.
\end{definition}

\begin{definition}[Decision for a low degree junta node]
\lean{FABL.lowDegreeJuntaNodeDecisionFromBatch}
\statementsource{Analysis-of-Boolean-Functions-by-Ryan-ODonnell}{
  pdf-ffbe9c7b3306-p081-b003,
  pdf-ffbe9c7b3306-p156-b006,
  pdf-ffbe9c7b3306-p156-b007,
  pdf-ffbe9c7b3306-p168-b001
}
\label{FABL.lowDegreeJuntaNodeDecisionFromBatch}
\leanok
\uses{FABL.JuntaFreeAssignment, FABL.JuntaNodeDecision, FABL.PositiveLearningParameter, FABL.Sign, FABL.binarySignEquiv, FABL.firstLowDegreeJuntaDetectedFrequency, FABL.firstNonconstantLowDegreeMonomial, FABL.juntaFreeCoordinateOfFin, FABL.lowDegreeJuntaPolynomialHypothesis, FABL.lowDegreeJuntaRestrictedSampleCount}
Let $n,k$ be nonnegative integers, let $P\subseteq\{0,\ldots,n-1\}$, let the failure
parameter lie in $(0,1/2]$, and take a batch of the prescribed restricted sample size of
free-coordinate sign assignments and labels. If the batch yields a detected frequency,
choose a free coordinate in its support. Otherwise form the low-degree polynomial
hypothesis from the batch: if it has a nonzero nonconstant monomial, choose a free
coordinate in that monomial's support; if it has none, return the constant sign given by
its coefficient at the empty set.
\end{definition}

\begin{definition}[Low degree junta polynomial work bound]
\lean{FABL.lowDegreeJuntaPolynomialWorkBound}
\statementsource{Analysis-of-Boolean-Functions-by-Ryan-ODonnell}{
  pdf-ffbe9c7b3306-p081-b003,
  pdf-ffbe9c7b3306-p156-b007,
  pdf-ffbe9c7b3306-p167-b005
}
\label{FABL.lowDegreeJuntaPolynomialWorkBound}
\leanok
\uses{FABL.PositiveLearningParameter, FABL.juntaRestrictionDimension, FABL.lowDegreeF₂MonomialCount, FABL.lowDegreeJuntaPolynomialDegree, FABL.lowDegreeJuntaPolynomialSampleCount}
Let $n,k$ be nonnegative integers, let $P\subseteq\{0,\ldots,n-1\}$, and let the failure
parameter lie in $(0,1/2]$. Write $r=|\{i\in\{0,\ldots,n-1\}:i\notin P\}|$,
$\ell=k-\lfloor3k/4\rfloor$, $m$ for the prescribed polynomial-learning sample count,
and $D$ for the number of subsets of an $r$-element set having size at most $\ell$. The
work bound is $m(D+1)+Dm(D+2)+D(D+1)$.
\end{definition}

\begin{definition}[Low degree junta node analysis work]
\lean{FABL.lowDegreeJuntaNodeAnalysisWork}
\statementsource{Analysis-of-Boolean-Functions-by-Ryan-ODonnell}{
  pdf-ffbe9c7b3306-p156-b006,
  pdf-ffbe9c7b3306-p168-b001
}
\label{FABL.lowDegreeJuntaNodeAnalysisWork}
\leanok
\uses{FABL.JuntaFreeAssignment, FABL.PositiveLearningParameter, FABL.Sign, FABL.juntaRestrictionDimension, FABL.lowDegreeF₂MonomialCount, FABL.lowDegreeF₂PolynomialLearnerWork, FABL.lowDegreeJuntaCoefficientSampleCount, FABL.lowDegreeJuntaCutoff, FABL.lowDegreeJuntaFourierFamily, FABL.lowDegreeJuntaPolynomialDegree, FABL.lowDegreeJuntaPolynomialSampleCount, FABL.lowDegreeJuntaPolynomialSamples, FABL.lowDegreeJuntaRestrictedSampleCount}
Let $n,k$ be nonnegative integers, let $P\subseteq\{0,\ldots,n-1\}$, let the failure
parameter lie in $(0,1/2]$, and take a batch of the prescribed restricted sample size.
Write $r$ for the number of coordinates outside $P$, $\ell=k-\lfloor3k/4\rfloor$, $F$
for the prescribed low-degree Fourier family, $c$ for the coefficient-estimation sample
count, $m$ for the polynomial-learning sample count, and $D$ for the number of subsets
of an $r$-element set of size at most $\ell$. The analysis work is
$|F|c(\lfloor3k/4\rfloor+1)$ plus the work of the low-degree polynomial learner on the
constructed polynomial samples, plus $|F|+D$.
\end{definition}

\begin{definition}[Low degree junta node analysis work bound]
\lean{FABL.lowDegreeJuntaNodeAnalysisWorkBound}
\statementsource{Analysis-of-Boolean-Functions-by-Ryan-ODonnell}{
  pdf-ffbe9c7b3306-p156-b006,
  pdf-ffbe9c7b3306-p168-b001
}
\label{FABL.lowDegreeJuntaNodeAnalysisWorkBound}
\leanok
\uses{FABL.PositiveLearningParameter, FABL.juntaRestrictionDimension, FABL.lowDegreeF₂MonomialCount, FABL.lowDegreeJuntaCoefficientSampleCount, FABL.lowDegreeJuntaCutoff, FABL.lowDegreeJuntaFourierFamily, FABL.lowDegreeJuntaPolynomialDegree, FABL.lowDegreeJuntaPolynomialWorkBound}
Let $n,k$ be nonnegative integers, let $P\subseteq\{0,\ldots,n-1\}$, and let the failure
parameter lie in $(0,1/2]$. Write $r$ for the number of coordinates outside $P$,
$\ell=k-\lfloor3k/4\rfloor$, $F$ for the prescribed low-degree Fourier family, $c$ for
the coefficient-estimation sample count, and $D$ for the number of subsets of an
$r$-element set of size at most $\ell$. The node-analysis work bound is
$|F|c(\lfloor3k/4\rfloor+1)$ plus the low-degree junta polynomial work bound, plus
$|F|+D$.
\end{definition}

\begin{theorem}[Node analysis work is bounded]
\lean{FABL.lowDegreeJuntaNodeAnalysisWork_le}
\label{FABL.lowDegreeJuntaNodeAnalysisWork_le}
\leanok
\uses{FABL.JuntaFreeAssignment, FABL.PositiveLearningParameter, FABL.Sign, FABL.juntaRestrictionDimension, FABL.lowDegreeF₂PolynomialLearnerWork_le, FABL.lowDegreeJuntaNodeAnalysisWork, FABL.lowDegreeJuntaNodeAnalysisWorkBound, FABL.lowDegreeJuntaPolynomialDegree, FABL.lowDegreeJuntaPolynomialSampleCount, FABL.lowDegreeJuntaPolynomialSamples, FABL.lowDegreeJuntaRestrictedSampleCount}
\difficulty{2}
Let $n,k$ be natural numbers, let $P$ be a set of coordinates among the $n$ coordinates,
let $\delta$ be a positive learning parameter with $0<\delta\le 1/2$, and let $B$ be a
batch of restricted examples of length equal to the restricted sample count for
$P,k,\delta$. The node analysis work on $B$ is at most the node analysis work bound for
$P,k,\delta$.
\end{theorem}

\begin{definition}[Continuation of junta node learning]
\lean{FABL.lowDegreeJuntaNodeContinuation}
\statementsource{Analysis-of-Boolean-Functions-by-Ryan-ODonnell}{
  pdf-ffbe9c7b3306-p156-b004,
  pdf-ffbe9c7b3306-p156-b005,
  pdf-ffbe9c7b3306-p168-b001
}
\label{FABL.lowDegreeJuntaNodeContinuation}
\leanok
\uses{FABL.JuntaFixedAssignment, FABL.JuntaNodeDecision, FABL.LearningProgram, FABL.MatchedJuntaExample, FABL.PositiveLearningParameter, FABL.lowDegreeJuntaNodeAnalysisWork, FABL.lowDegreeJuntaNodeDecisionFromBatch, FABL.lowDegreeJuntaRestrictedSampleCount, FABL.matchedJuntaRestrictionExample}
Let $n,k$ be natural numbers, let $P$ be a set of coordinates in $\text{Fin }n$, let $z$
assign a sign to each coordinate in $P$, and let the failure parameter lie in $(0,1/2]$.
For an optional batch of the prescribed restricted sample size, whose examples match $z$
on $P$ when the batch is present, this definition returns the empty decision in a
random-example learning program if no batch is supplied. If a batch is supplied, it
forms the restricted examples, charges the node-analysis work, and returns the
junta-node decision computed from that batch.
\end{definition}

\begin{definition}[Junta relevant-coordinate program]
\lean{FABL.lowDegreeJuntaRelevantCoordinateProgram}
\statementsource{Analysis-of-Boolean-Functions-by-Ryan-ODonnell}{
  pdf-ffbe9c7b3306-p156-b005,
  pdf-ffbe9c7b3306-p156-b006,
  pdf-ffbe9c7b3306-p156-b007,
  pdf-ffbe9c7b3306-p156-b008,
  pdf-ffbe9c7b3306-p168-b001
}
\label{FABL.lowDegreeJuntaRelevantCoordinateProgram}
\leanok
\uses{FABL.JuntaFixedAssignment, FABL.JuntaNodeDecision, FABL.LearningProgram, FABL.LearningProgram.bind, FABL.PositiveLearningParameter, FABL.juntaRestrictionSampleProgram, FABL.lowDegreeJuntaHalfFailure, FABL.lowDegreeJuntaNodeContinuation, FABL.lowDegreeJuntaRestrictedSampleCount}
Let $n,k$ be nonnegative integers, let $P\subseteq\{0,\ldots,n-1\}$ be a set of
coordinates, let $z$ assign a sign to each coordinate in $P$, and let $\delta\in(0,1/2]$
be a failure parameter. The program returns either no decision or a decision that the
restriction is constant or that it splits on a coordinate outside $P$; it is defined by
binding the low-degree junta node continuation to the junta restriction sample program,
using the restricted sample count and half the failure parameter.
\end{definition}

\begin{definition}[Junta node sample count]
\lean{FABL.lowDegreeJuntaNodeRandomExamples}
\statementsource{Analysis-of-Boolean-Functions-by-Ryan-ODonnell}{
  pdf-ffbe9c7b3306-p081-b003,
  pdf-ffbe9c7b3306-p156-b005,
  pdf-ffbe9c7b3306-p156-b006,
  pdf-ffbe9c7b3306-p168-b001
}
\label{FABL.lowDegreeJuntaNodeRandomExamples}
\leanok
\uses{FABL.PositiveLearningParameter, FABL.juntaRestrictionSampleCount, FABL.lowDegreeJuntaHalfFailure, FABL.lowDegreeJuntaRestrictedSampleCount}
Let $n,k$ be nonnegative integers, let $P\subseteq\{0,\ldots,n-1\}$ be a set of
coordinates, and let $\delta\in(0,1/2]$ be a failure parameter. The junta node
random-example count is the junta restriction sample count at degree parameter $k$, with
the restricted sample count for $P,k,\delta$ and half the failure parameter as inputs.
\end{definition}

\begin{definition}[Junta node work bound]
\lean{FABL.lowDegreeJuntaNodeWorkBound}
\statementsource{Analysis-of-Boolean-Functions-by-Ryan-ODonnell}{
  pdf-ffbe9c7b3306-p156-b006,
  pdf-ffbe9c7b3306-p156-b007,
  pdf-ffbe9c7b3306-p156-b008,
  pdf-ffbe9c7b3306-p168-b001
}
\label{FABL.lowDegreeJuntaNodeWorkBound}
\leanok
\uses{FABL.PositiveLearningParameter, FABL.juntaRestrictionSampleCount, FABL.juntaRestrictionSampleWork, FABL.lowDegreeJuntaHalfFailure, FABL.lowDegreeJuntaNodeAnalysisWorkBound, FABL.lowDegreeJuntaRestrictedSampleCount}
Let $n,k$ be nonnegative integers, let $P\subseteq\{0,\ldots,n-1\}$ be a set of
coordinates, and let $\delta\in(0,1/2]$ be a failure parameter. Set $M$ equal to the
restricted sample count for $P,k,\delta$. The node work bound is the sum of the junta
restriction sample count at degree parameter $k$ with inputs $M$ and half the failure
parameter, the corresponding junta restriction sample-work bound, and the node analysis
work bound for $P,k,\delta$.
\end{definition}

\begin{definition}[Uniform junta node sample bound]
\lean{FABL.lowDegreeJuntaNodeRandomExampleBound}
\statementsource{Analysis-of-Boolean-Functions-by-Ryan-ODonnell}{
  pdf-ffbe9c7b3306-p081-b003,
  pdf-ffbe9c7b3306-p156-b005,
  pdf-ffbe9c7b3306-p168-b001
}
\label{FABL.lowDegreeJuntaNodeRandomExampleBound}
\leanok
\uses{FABL.PositiveLearningParameter, FABL.lowDegreeJuntaNodeRandomExamples}
Let $n,k$ be nonnegative integers and let $\delta\in(0,1/2]$ be a failure parameter. The
uniform node random-example bound is the maximum, over all sets
$P\subseteq\{0,\ldots,n-1\}$, of the junta node random-example count for $P,k,\delta$.
\end{definition}

\begin{definition}[Uniform junta node work bound]
\lean{FABL.lowDegreeJuntaNodeUniformWorkBound}
\statementsource{Analysis-of-Boolean-Functions-by-Ryan-ODonnell}{pdf-ffbe9c7b3306-p168-b001}
\label{FABL.lowDegreeJuntaNodeUniformWorkBound}
\leanok
\uses{FABL.PositiveLearningParameter, FABL.lowDegreeJuntaNodeWorkBound}
Let $n,k$ be nonnegative integers and let $\delta\in(0,1/2]$ be a failure parameter. The
uniform node work bound is the maximum, over all sets $P\subseteq\{0,\ldots,n-1\}$, of
the junta node work bound for $P,k,\delta$.
\end{definition}

\begin{definition}[Reindexed junta witness coordinates]
\lean{FABL.reindexedJuntaWitnessCoordinates}
\statementsource{Analysis-of-Boolean-Functions-by-Ryan-ODonnell}{pdf-ffbe9c7b3306-p168-b001}
\label{FABL.reindexedJuntaWitnessCoordinates}
\leanok
\uses{FABL.FixedIndex, FABL.juntaRestrictionDimension}
Let $n$ be a nonnegative integer, and let $P,J\subseteq\{0,\ldots,n-1\}$ be coordinate
sets. Enumerate the coordinates outside $P$ by positions in a set of size equal to the
number of coordinates outside $P$. The reindexed witness coordinates are the positions
in this enumeration whose corresponding original coordinates belong to $J$.
\end{definition}

\begin{theorem}[Bound on reindexed witness coordinates]
\lean{FABL.card_reindexedJuntaWitnessCoordinates_le}
\label{FABL.card_reindexedJuntaWitnessCoordinates_le}
\leanok
\uses{FABL.FixedIndex, FABL.juntaRestrictionDimension, FABL.reindexedJuntaWitnessCoordinates}
\difficulty{3}
Let $n$ be a natural number and let $P,J\subseteq\{0,\ldots,n-1\}$. Reindex the
coordinates outside $P$ by their positions in a set of size equal to the number of
coordinates outside $P$, and select those whose original indices lie in $J$. The number
selected is at most $|J|$.
\end{theorem}

\begin{theorem}[Junta Property of a Reindexed Restriction]
\lean{FABL.reindexedJuntaRestriction_isKJunta}
\label{FABL.reindexedJuntaRestriction_isKJunta}
\leanok
\uses{FABL.BooleanFunction, FABL.FixedIndex, FABL.IsKJunta, FABL.JuntaFixedAssignment, FABL.card_reindexedJuntaWitnessCoordinates_le, FABL.combineJuntaAssignment, FABL.fixedSignCubeEquiv, FABL.juntaRestrictionDimension, FABL.reindexedJuntaRestriction, FABL.reindexedJuntaWitnessCoordinates}
\difficulty{4}
Let $n$ be a natural number, let $f\colon\{-1,1\}^n\to\{-1,1\}$, let $J$ and $P$ be sets
of coordinates, let $z$ assign a sign to every coordinate in $P$, and let $k$ be a
natural number. If $|J|\le k$ and $f$ depends only on the coordinates in $J$, then the
restriction of $f$ obtained by fixing the coordinates in $P$ to $z$ and reindexing the
remaining coordinates is a $k$-junta.
\end{theorem}

\begin{theorem}[Relevance of a Junta Restriction]
\lean{FABL.isRelevantJuntaRestriction_of_reindexed}
\label{FABL.isRelevantJuntaRestriction_of_reindexed}
\leanok
\uses{FABL.BooleanFunction, FABL.BooleanFunction.toReal, FABL.FixedIndex, FABL.IsRelevant, FABL.IsRelevantJuntaRestriction, FABL.JuntaFixedAssignment, FABL.JuntaFreeAssignment, FABL.fixedSignCubeEquiv, FABL.isRelevant_iff_exists_setCoordinate_ne, FABL.juntaFreeCoordinateOfFin, FABL.juntaRestriction, FABL.juntaRestrictionDimension, FABL.reindexedJuntaRestriction, FABL.setCoordinate, FABL.signValue}
\difficulty{4}
Let $n$ be a natural number, let $f$ be a Boolean function on the sign cube
$\\{−1,1\\}^n$, let $P\\subseteq\\{1,\\ldots,n\\}$ be a set of fixed coordinates, and
fix signs on $P$. Reindex the restriction of $f$ to the remaining coordinates as a
function on a sign cube of dimension equal to the number of remaining coordinates. If
this reindexed function has positive influence in coordinate $i$, then the restricted
function depends on the corresponding free coordinate: there are two assignments to the
free coordinates that differ only there and on which the restricted function takes
different values.
\end{theorem}

\begin{definition}[Accuracy of low-degree junta Fourier estimates]
\lean{FABL.LowDegreeJuntaFourierEstimatesAccurate}
\statementsource{Analysis-of-Boolean-Functions-by-Ryan-ODonnell}{
  pdf-ffbe9c7b3306-p080-b002,
  pdf-ffbe9c7b3306-p156-b006,
  pdf-ffbe9c7b3306-p167-b005
}
\label{FABL.LowDegreeJuntaFourierEstimatesAccurate}
\leanok
\uses{FABL.BooleanFunction, FABL.BooleanFunction.toReal, FABL.JuntaFixedAssignment, FABL.JuntaFreeAssignment, FABL.PositiveLearningParameter, FABL.Sign, FABL.fourierCoeff, FABL.lowDegreeJuntaCoefficientAccuracy, FABL.lowDegreeJuntaCoefficientEstimate, FABL.lowDegreeJuntaFourierFamily, FABL.lowDegreeJuntaRestrictedSampleCount, FABL.reindexedJuntaRestriction}
Let $n,k$ be nonnegative integers, let $f:\{−1,1\}^n\to\{−1,1\}$ be a Boolean function,
let $P\subseteq\{0,\ldots,n-1\}$ be a coordinate set, let $z$ assign signs on $P$, and
let $\delta\in(0,1/2]$. For every batch of the specified restricted-sample size
consisting of assignments on coordinates outside $P$ paired with signs, require that
every Fourier coefficient estimate indexed by a set in the low-degree junta Fourier
family differs from the corresponding Fourier coefficient of the restriction of $f$
under $z$ by less than the coefficient accuracy parameter.
\end{definition}

\begin{definition}[Exact low-degree junta polynomial hypothesis]
\lean{FABL.LowDegreeJuntaPolynomialHypothesisExact}
\statementsource{Analysis-of-Boolean-Functions-by-Ryan-ODonnell}{
  pdf-ffbe9c7b3306-p080-b004,
  pdf-ffbe9c7b3306-p081-b003,
  pdf-ffbe9c7b3306-p156-b007,
  pdf-ffbe9c7b3306-p167-b005
}
\label{FABL.LowDegreeJuntaPolynomialHypothesisExact}
\leanok
\uses{FABL.BooleanFunction, FABL.JuntaFixedAssignment, FABL.JuntaFreeAssignment, FABL.LowDegreeF₂Hypothesis.evaluate, FABL.PositiveLearningParameter, FABL.Sign, FABL.booleanFunctionF₂Encoding, FABL.lowDegreeJuntaPolynomialHypothesis, FABL.lowDegreeJuntaRestrictedSampleCount, FABL.reindexedJuntaRestriction}
Let $n,k$ be nonnegative integers, let $f:\{−1,1\}^n\to\{−1,1\}$ be a Boolean function,
let $P\subseteq\{0,\ldots,n-1\}$ be a coordinate set, let $z$ assign signs on $P$, and
let $\delta\in(0,1/2]$. For every batch of the specified restricted-sample size
consisting of assignments on coordinates outside $P$ paired with signs, the polynomial
hypothesis constructed from the batch evaluates exactly to the Boolean encoding of the
restriction of $f$ under $z$.
\end{definition}

\begin{theorem}[Junta Fourier coefficient estimation]
\lean{FABL.measure_lowDegreeJuntaCoefficientEstimate_bad_le}
\proofsource{Analysis-of-Boolean-Functions-by-Ryan-ODonnell}{
  pdf-ffbe9c7b3306-p080-b003,
  pdf-ffbe9c7b3306-p156-b006
}
\label{FABL.measure_lowDegreeJuntaCoefficientEstimate_bad_le}
\leanok
\uses{FABL.BooleanFunction, FABL.BooleanFunction.toReal, FABL.JuntaFixedAssignment, FABL.JuntaFreeAssignment, FABL.PositiveLearningParameter, FABL.fourierCoeff, FABL.fourierEstimatorSampleCount_pos, FABL.juntaRestrictionDimension, FABL.juntaRestrictionLabeledBatch, FABL.lowDegreeJuntaCoefficientAccuracy, FABL.lowDegreeJuntaCoefficientConfidence, FABL.lowDegreeJuntaCoefficientEstimate, FABL.lowDegreeJuntaCoefficientEstimate_cast_ideal, FABL.lowDegreeJuntaCoefficientInputRow, FABL.lowDegreeJuntaCoefficientSampleCount, FABL.lowDegreeJuntaFourierFamily, FABL.lowDegreeJuntaRestrictedSampleCount, FABL.map_uniformPMF_lowDegreeJuntaCoefficientInputRow, FABL.measure_realEmpiricalFourierCoeff_sub_ge_le, FABL.positiveLearningParameter_toReal_mem_Ioc, FABL.realEmpiricalFourierCoeff, FABL.reindexedJuntaRestriction, FABL.two_mul_exp_neg_fourierEstimatorSampleCount_le, FABL.uniformPMF}
\difficulty{4}
Let $n\in\mathbb{N}$, let $f:\{-1,1\}^n\to\{-1,1\}$, let $P\subseteq[n]$, and fix signs
on the coordinates in $P$. Let $k\in\mathbb{N}$, let the rational failure parameter
$\eta$ satisfy $0<\eta\le 1/2$, and choose any Fourier index $S$ in the low-degree
family on the remaining coordinates, whose members have size at most $\lfloor
3k/4\rfloor$. For the uniform batch of free-coordinate inputs of the sample size
specified by the low-degree junta coefficient estimator, the probability that its
empirical estimate of the $S$-th Fourier coefficient of the restricted function differs
from the true coefficient by at least $1/(3\cdot 2^k)$ is at most the estimator's
confidence parameter (the quarter-failure parameter divided by the size of that Fourier
family).
\end{theorem}

\begin{theorem}[Fourier estimate failure probability]
\lean{FABL.idealLowDegreeJuntaFourierInaccuracyProbability_le}
\proofsource{Analysis-of-Boolean-Functions-by-Ryan-ODonnell}{
  pdf-ffbe9c7b3306-p156-b006,
  pdf-ffbe9c7b3306-p080-b003
}
\label{FABL.idealLowDegreeJuntaFourierInaccuracyProbability_le}
\leanok
\uses{FABL.BooleanFunction, FABL.BooleanFunction.toReal, FABL.JuntaFixedAssignment, FABL.JuntaFreeAssignment, FABL.LowDegreeJuntaFourierEstimatesAccurate, FABL.PositiveLearningParameter, FABL.fourierCoeff, FABL.juntaRestrictionLabeledBatch, FABL.lowDegreeJuntaCoefficientAccuracy, FABL.lowDegreeJuntaCoefficientConfidence, FABL.lowDegreeJuntaCoefficientConfidence_value, FABL.lowDegreeJuntaCoefficientEstimate, FABL.lowDegreeJuntaFourierFamily, FABL.lowDegreeJuntaFourierFamily_nonempty, FABL.lowDegreeJuntaQuarterFailure, FABL.lowDegreeJuntaQuarterFailure_value, FABL.lowDegreeJuntaRestrictedSampleCount, FABL.measure_lowDegreeJuntaCoefficientEstimate_bad_le, FABL.reindexedJuntaRestriction, FABL.uniformPMF}
\difficulty{5}
Let $n\ge 0$, let $f\colon\{\pm1\}^n\to\{\pm1\}$, let $P\subseteq\{1,\ldots,n\}$ be a
set of coordinates, let $z\colon P\to\{\pm1\}$ be a fixed assignment, let $k\ge0$, and
let $\delta$ be a rational number with $0<\delta\le1/2$. Draw uniformly a batch of the
prescribed size whose inputs assign signs to the coordinates outside $P$, and label each
input by the restriction of $f$ with the coordinates in $P$ fixed according to $z$. The
probability that at least one Fourier coefficient in the low-degree family on the
remaining coordinates is estimated with error at least the prescribed accuracy is at
most $\delta/4$.
\end{theorem}

\begin{theorem}[Polynomial learner failure probability]
\lean{FABL.idealLowDegreeJuntaPolynomialFailureProbability_le}
\label{FABL.idealLowDegreeJuntaPolynomialFailureProbability_le}
\leanok
\uses{FABL.BooleanFunction, FABL.JuntaFixedAssignment, FABL.JuntaFreeAssignment, FABL.LowDegreeF₂Hypothesis.evaluate, FABL.LowDegreeJuntaPolynomialHypothesisExact, FABL.PositiveLearningParameter, FABL.booleanFunctionF₂Encoding, FABL.functionAlgebraicDegree, FABL.juntaRestrictionDimension, FABL.juntaRestrictionLabeledBatch, FABL.lowDegreeF₂LearningSampleCount, FABL.lowDegreeF₂PolynomialLearnerLabeledOutput, FABL.lowDegreeF₂PolynomialLearner_failure_subset_separationFailure, FABL.lowDegreeJuntaPolynomialDegree, FABL.lowDegreeJuntaPolynomialHypothesis_ideal, FABL.lowDegreeJuntaPolynomialInputRow, FABL.lowDegreeJuntaPolynomialSampleCount, FABL.lowDegreeJuntaQuarterFailure, FABL.lowDegreeJuntaRestrictedSampleCount, FABL.map_uniformPMF_lowDegreeJuntaPolynomialInputRow, FABL.measure_signSampleLowDegreeF₂SeparationFailureSet_scheduled_le, FABL.reindexedJuntaRestriction, FABL.signSampleLowDegreeF₂SeparationFailureSet, FABL.uniformPMF}
\difficulty{5}
Let $n\ge0$, let $f\colon\{\pm1\}^n\to\{\pm1\}$, let $P\subseteq\{1,\ldots,n\}$, let
$z\colon P\to\{\pm1\}$, and let $k\ge0$. Suppose the algebraic degree over
$\mathbb{F}_2$ of the Boolean function obtained by restricting $f$ according to $z$ and
encoding it on the remaining coordinates is at most the prescribed polynomial degree.
For any rational failure parameter $\delta$ with $0<\delta\le1/2$, the probability that
the low-degree polynomial learner fails to recover that restricted function exactly,
when supplied the prescribed uniformly sampled batch, is at most $\delta/4$.
\end{theorem}

\begin{theorem}[Detected frequencies have nonzero coefficients]
\lean{FABL.fourierCoeff_ne_zero_of_mem_lowDegreeJuntaDetectedFrequencies}
\proofsource{Analysis-of-Boolean-Functions-by-Ryan-ODonnell}{pdf-ffbe9c7b3306-p156-b006}
\label{FABL.fourierCoeff_ne_zero_of_mem_lowDegreeJuntaDetectedFrequencies}
\leanok
\uses{FABL.BooleanFunction, FABL.BooleanFunction.toReal, FABL.JuntaFixedAssignment, FABL.JuntaFreeAssignment, FABL.LowDegreeJuntaFourierEstimatesAccurate, FABL.PositiveLearningParameter, FABL.Sign, FABL.fourierCoeff, FABL.juntaRestrictionDimension, FABL.lowDegreeJuntaAccuracy_lt_detectionThreshold, FABL.lowDegreeJuntaCoefficientAccuracy, FABL.lowDegreeJuntaCoefficientEstimate, FABL.lowDegreeJuntaDetectedFrequencies, FABL.lowDegreeJuntaDetectionThreshold, FABL.lowDegreeJuntaFourierFamily, FABL.lowDegreeJuntaRestrictedSampleCount, FABL.reindexedJuntaRestriction}
\difficulty{3}
Let $n,k$ be nonnegative integers, let $f$ be a sign-valued function on the
$n$-dimensional sign cube, let $P$ be a set of coordinates, and fix a sign assignment
$z$ on $P$. Let $\delta$ be a rational number with $0<\delta\le 1/2$, and take a sample
batch of the prescribed size for the restriction of $f$ obtained by fixing the
coordinates in $P$ according to $z$. Suppose that, for every coordinate set of size at
most $\lfloor 3k/4\rfloor$, its batch-based Fourier estimate differs from the true
Fourier coefficient of this restricted function by less than $1/(3\cdot 2^k)$. Then
every nonempty coordinate set of size at most $\lfloor 3k/4\rfloor$ whose estimated
coefficient has absolute value at least $1/(2\cdot 2^k)$ has a nonzero true Fourier
coefficient for the restricted function.
\end{theorem}

\begin{theorem}[Undetected low-degree coefficients vanish]
\lean{FABL.fourierCoeff_eq_zero_of_no_lowDegreeJuntaDetection}
\proofsource{Analysis-of-Boolean-Functions-by-Ryan-ODonnell}{
  pdf-ffbe9c7b3306-p156-b006,
  pdf-ffbe9c7b3306-p156-b007
}
\label{FABL.fourierCoeff_eq_zero_of_no_lowDegreeJuntaDetection}
\leanok
\uses{FABL.BooleanFunction, FABL.BooleanFunction.toReal, FABL.IsKJunta, FABL.JuntaFixedAssignment, FABL.JuntaFreeAssignment, FABL.LowDegreeJuntaFourierEstimatesAccurate, FABL.PositiveLearningParameter, FABL.Sign, FABL.firstLowDegreeJuntaDetectedFrequency, FABL.fourierCoeff, FABL.fourierCoeff_eq_zero_or_inv_two_pow_le_abs_of_isKJunta, FABL.juntaRestrictionDimension, FABL.lowDegreeJuntaCoefficientAccuracy, FABL.lowDegreeJuntaCoefficientEstimate, FABL.lowDegreeJuntaCutoff, FABL.lowDegreeJuntaDetectedFrequencies, FABL.lowDegreeJuntaDetectionThreshold, FABL.lowDegreeJuntaFourierFamily, FABL.lowDegreeJuntaRestrictedSampleCount, FABL.mem_lowDegreeJuntaFourierFamily, FABL.reindexedJuntaRestriction, FABL.reindexedJuntaRestriction_isKJunta}
\difficulty{6}
Let $n\ge 0$, let the target be a sign-valued function on the $n$-dimensional sign cube,
and let $J$ and $P$ be coordinate sets. Fix signs on $P$, let $k$ be a natural number,
and let the failure parameter lie strictly between $0$ and $1/2$. Suppose $|J|\le k$ and
the target depends only on coordinates in $J$. For a batch of the prescribed finite
number of examples on the coordinates outside $P$, assume that every estimated Fourier
coefficient in the low-degree family for the resulting restricted target differs from
its true coefficient by less than the prescribed accuracy, and that the procedure for
detecting a nonempty frequency returns none. Then every nonempty coordinate set $S$ of
size at most $\lfloor 3k/4\rfloor$ has zero Fourier coefficient for that restricted
target, viewed as real-valued.
\end{theorem}

\begin{theorem}[Correlation immunity from no detected frequency]
\lean{FABL.isCorrelationImmune_of_no_lowDegreeJuntaDetection}
\label{FABL.isCorrelationImmune_of_no_lowDegreeJuntaDetection}
\leanok
\uses{FABL.BooleanFunction, FABL.IsCorrelationImmune, FABL.JuntaFixedAssignment, FABL.JuntaFreeAssignment, FABL.LowDegreeJuntaFourierEstimatesAccurate, FABL.PositiveLearningParameter, FABL.Sign, FABL.firstLowDegreeJuntaDetectedFrequency, FABL.fourierCoeff_eq_zero_of_no_lowDegreeJuntaDetection, FABL.lowDegreeJuntaCutoff, FABL.lowDegreeJuntaRestrictedSampleCount, FABL.reindexedJuntaRestriction}
\difficulty{3}
Let $n,k$ be natural numbers, let the target be a sign-valued function on $\{-1,1\}^n$,
and let $J,P\subseteq\{1,\ldots,n\}$ be coordinate sets. Fix signs on $P$, assume
$|J|\le k$, and assume the target depends only on coordinates in $J$. Given a batch of
samples indexed by the prescribed sample count, suppose every Fourier estimate in the
prescribed low-degree family for the restriction obtained by fixing those signs on $P$
has error strictly less than the prescribed accuracy, and the frequency detector returns
none. Then every nonempty Fourier coefficient of that restricted function on a set of
size at most $\lfloor 3k/4\rfloor$ is zero.
\end{theorem}

\begin{definition}[Embedding of witness coordinates]
\lean{FABL.juntaWitnessCoordinateEmbedding}
\label{FABL.juntaWitnessCoordinateEmbedding}
\leanok
Let $r$ be a natural number and let $J$ be a set of coordinates in $0,\\ldots,r-1$.
Enumerating the elements of $J$ gives an embedding of the set of indices from zero to
$|J|-1$ into the coordinates from zero to $r-1$, with image exactly $J$.
\end{definition}

\begin{definition}[Restricted Boolean function]
\lean{FABL.juntaWitnessCore}
\statementsource{Analysis-of-Boolean-Functions-by-Ryan-ODonnell}{
  pdf-ffbe9c7b3306-p150-b004,
  pdf-ffbe9c7b3306-p150-b005,
  pdf-ffbe9c7b3306-p170-b004
}
\label{FABL.juntaWitnessCore}
\leanok
\uses{FABL.BooleanFunction, FABL.signRestriction}
Let $r$ be a natural number, let $g$ be a Boolean function on the $r$-dimensional sign
cube, and let $J$ be a set of coordinates in that cube. The resulting Boolean function
on a cube with $|J|$ coordinates is obtained by evaluating $g$ with every coordinate in
$J$ fixed to $+1$, and with the remaining coordinates supplied by the input, reindexed
by their positions in $J$.
\end{definition}

\begin{theorem}[Degree under coordinate pullback]
\lean{FABL.functionAlgebraicDegree_coordinatePullback_le}
\label{FABL.functionAlgebraicDegree_coordinatePullback_le}
\leanok
\uses{FABL.F₂BooleanFunction, FABL.F₂Cube, FABL.algebraicDegree_le_iff, FABL.anfCoeff, FABL.anfEval, FABL.anfEval_anfCoeff, FABL.anfMonomial, FABL.functionAlgebraicDegree, FABL.functionAlgebraicDegree_anfMonomial, FABL.functionAlgebraicDegree_finset_sum_le, FABL.functionAlgebraicDegree_zero}
\difficulty{5}
Let $r,s\ge 0$, let $f:\mathbb{F}_2^s\to\mathbb{F}_2$ be a Boolean function, and let
$e:\{1,\ldots,s\}\hookrightarrow\{1,\ldots,r\}$ be an injection. The algebraic degree
over $\mathbb{F}_2$ of the function on $\mathbb{F}_2^r$ obtained by evaluating $f$ on
the coordinates selected by $e$ is at most the algebraic degree of $f$.
\end{theorem}

\begin{theorem}[Encoding a junta core]
\lean{FABL.booleanFunctionF₂Encoding_eq_juntaWitnessCore_pullback}
\leanok
\uses{FABL.BooleanFunction, FABL.F₂Cube, FABL.binaryCubeSignEquiv, FABL.binaryCubeSignEquiv_apply, FABL.binarySignEquiv, FABL.booleanFunctionF₂Encoding, FABL.combineSignCube, FABL.combineSignCube_apply_free, FABL.juntaWitnessCoordinateEmbedding, FABL.juntaWitnessCore}
\difficulty{4}
Let $r$ be a nonnegative integer, let $J$ be a subset of the coordinates
$\{0,\ldots,r-1\}$, and let $g$ be a sign-valued Boolean function on the $r$-dimensional
sign cube that depends only on coordinates in $J$. The binary encoding of $g$ equals the
function that first restricts its input to the coordinates indexed by $J$, in their
increasing order, and then evaluates the binary encoding of the corresponding core of
$g$ on the $|J|$-dimensional sign cube.
\end{theorem}

\begin{theorem}[Degree and a junta witness core]
\lean{FABL.functionAlgebraicDegree_encoding_le_juntaWitnessCore}
\label{FABL.functionAlgebraicDegree_encoding_le_juntaWitnessCore}
\leanok
\uses{FABL.BooleanFunction, FABL.booleanFunctionF₂Encoding, FABL.booleanFunctionF₂Encoding_eq_juntaWitnessCore_pullback, FABL.functionAlgebraicDegree, FABL.functionAlgebraicDegree_coordinatePullback_le, FABL.juntaWitnessCoordinateEmbedding, FABL.juntaWitnessCore}
\difficulty{1}
Let $r\ge 0$, let $g:\{\pm1\}^r\to\{\pm1\}$, and let $J$ be a set of coordinates in
$\{1,\ldots,r\}$. Suppose $g$ depends only on coordinates in $J$. Fix every coordinate
outside $J$ to $+1$ and regard the resulting function as a function on the coordinates
in $J$. The algebraic degree over $\mathbb{F}_2$ of the encoding of $g$ is at most that
of the encoding of this restricted function.
\end{theorem}

\begin{theorem}[Fourier coefficients of a junta witness]
\lean{FABL.fourierCoeff_juntaWitnessCore}
\label{FABL.fourierCoeff_juntaWitnessCore}
\leanok
\uses{FABL.BooleanFunction, FABL.BooleanFunction.toReal, FABL.FixedSignCube, FABL.FreeSignCube, FABL.Sign, FABL.combineSignCube, FABL.combineSignCube_apply_free, FABL.expect_restrictionFourierCoeff, FABL.fourierCoeff, FABL.indexedFourierCoeff, FABL.indexedMonomial, FABL.juntaWitnessCoordinateEmbedding, FABL.juntaWitnessCore, FABL.liftFreeFrequency, FABL.monomial, FABL.restrictionFourierCoeff, FABL.signRestriction, FABL.signValue}
\difficulty{6}
Let $r$ be a nonnegative integer, let $g:\{\pm1\}^r\to\{\pm1\}$, and let
$J\subseteq\{0,\ldots,r-1\}$. Assume that $g$ depends only on coordinates in $J$. For
every subset $T$ of the $|J|$ coordinates of the function obtained by restricting $g$ to
$J$ with all coordinates outside $J$ fixed to $1$, its Fourier coefficient at $T$ equals
the Fourier coefficient of $g$ at the corresponding subset of coordinates of $J$, using
the increasing enumeration of $J$.
\end{theorem}

\begin{theorem}[Correlation immunity of the junta witness core]
\lean{FABL.juntaWitnessCore_isCorrelationImmune}
\label{FABL.juntaWitnessCore_isCorrelationImmune}
\leanok
\uses{FABL.BooleanFunction, FABL.IsCorrelationImmune, FABL.fourierCoeff_juntaWitnessCore, FABL.juntaWitnessCoordinateEmbedding, FABL.juntaWitnessCore}
\difficulty{4}
Let $r$ and $d$ be nonnegative integers, let $g$ be a Boolean function on the
$r$-dimensional sign cube, and let $J$ be a set of coordinates. Suppose that $g$ depends
only on coordinates in $J$ and is correlation immune through degree $d$, meaning that
every nonempty Fourier coefficient indexed by a set of at most $d$ coordinates is zero.
Then the function on the $|J|$-dimensional sign cube obtained by restricting the
coordinates outside $J$ to $+1$ and reindexing the coordinates in $J$ is also
correlation immune through degree $d$.
\end{theorem}

\begin{theorem}[Degree of a constant Boolean function]
\lean{FABL.functionAlgebraicDegree_encoding_eq_zero_of_eq_const}
\label{FABL.functionAlgebraicDegree_encoding_eq_zero_of_eq_const}
\leanok
\uses{FABL.BooleanFunction, FABL.F₂Cube, FABL.Sign, FABL.binarySignEquiv, FABL.booleanFunctionF₂Encoding, FABL.functionAlgebraicDegree, FABL.functionAlgebraicDegree_one, FABL.functionAlgebraicDegree_zero}
\difficulty{4}
Let $r\ge 0$, and let $g:\{\pm1\}^r\to\{\pm1\}$ be a Boolean function that is constantly
equal to some $v\in\{\pm1\}$. Its encoding as a function from $\mathbb{F}_2^r$ to
$\mathbb{F}_2$ has algebraic degree $0$.
\end{theorem}

\begin{theorem}[Degree bound for correlation-immune juntas]
\lean{FABL.functionAlgebraicDegree_encoding_le_sub_of_isKJunta_of_isCorrelationImmune}
\proofsource{Analysis-of-Boolean-Functions-by-Ryan-ODonnell}{
  pdf-ffbe9c7b3306-p150-b004,
  pdf-ffbe9c7b3306-p150-b005
}
\label{FABL.functionAlgebraicDegree_encoding_le_sub_of_isKJunta_of_isCorrelationImmune}
\leanok
\uses{FABL.BooleanFunction, FABL.BooleanFunction.toReal, FABL.IsCorrelationImmune, FABL.IsFourierRegular, FABL.IsKJunta, FABL.booleanFunctionF₂Encoding, FABL.functionAlgebraicDegree, FABL.functionAlgebraicDegree_booleanFunctionF₂Encoding_le_of_isCorrelationImmune, FABL.functionAlgebraicDegree_encoding_eq_zero_of_eq_const, FABL.functionAlgebraicDegree_encoding_le_juntaWitnessCore, FABL.isFourierRegular_zero_iff_exists_const, FABL.isLowDegreeFourierRegular_dimension_iff, FABL.juntaWitnessCore, FABL.juntaWitnessCore_isCorrelationImmune, FABL.signValue_injective}
\difficulty{5}
Let $r,k,d\ge 0$, and let $g:\{\pm1\}^r\to\{\pm1\}$. Suppose $g$ depends on at most $k$
coordinates. Also suppose that every nonempty set of at most $d$ coordinates has Fourier
coefficient zero for the real-valued function obtained from $g$ by interpreting its
outputs as signs. Then the algebraic degree over $\mathbb{F}_2$ of the encoding of $g$
is at most $k-d$, with subtraction taken in the natural numbers.
\end{theorem}

\begin{theorem}[Function algebraic degree reindexed junta restriction le of no detection]
\lean{FABL.functionAlgebraicDegree_reindexedJuntaRestriction_le_of_no_detection}
\statementsource{Analysis-of-Boolean-Functions-by-Ryan-ODonnell}{
  pdf-ffbe9c7b3306-p156-b007,
  pdf-ffbe9c7b3306-p170-b002
}
\label{FABL.functionAlgebraicDegree_reindexedJuntaRestriction_le_of_no_detection}
\leanok
\uses{FABL.BooleanFunction, FABL.JuntaFixedAssignment, FABL.JuntaFreeAssignment, FABL.LowDegreeJuntaFourierEstimatesAccurate, FABL.PositiveLearningParameter, FABL.Sign, FABL.booleanFunctionF₂Encoding, FABL.firstLowDegreeJuntaDetectedFrequency, FABL.functionAlgebraicDegree, FABL.functionAlgebraicDegree_encoding_le_sub_of_isKJunta_of_isCorrelationImmune, FABL.isCorrelationImmune_of_no_lowDegreeJuntaDetection, FABL.lowDegreeJuntaPolynomialDegree, FABL.lowDegreeJuntaRestrictedSampleCount, FABL.reindexedJuntaRestriction, FABL.reindexedJuntaRestriction_isKJunta}
\difficulty{1}
% dataset-pipeline: stub -- prose pending build_informal_statements.py
Statement pending regeneration.
\end{theorem}

\begin{theorem}[A detected frequency selected by the first-detection rule]
\lean{FABL.mem_detectedFrequencies_of_firstDetected_eq_some}
\label{FABL.mem_detectedFrequencies_of_firstDetected_eq_some}
\leanok
\uses{FABL.JuntaFreeAssignment, FABL.PositiveLearningParameter, FABL.Sign, FABL.firstLowDegreeJuntaDetectedFrequency, FABL.juntaRestrictionDimension, FABL.lowDegreeJuntaDetectedFrequencies, FABL.lowDegreeJuntaRestrictedSampleCount}
\difficulty{5}
Let $P$ be a subset of $\{0,\ldots,n-1\}$, let $k$ be a natural number, let the failure
parameter lie in $(0,1/2]$, and let a batch of the prescribed size of assignments on the
coordinates outside $P$, paired with signs, be given. If the first detected frequency
selected from this batch is $S$, then $S$ belongs to the batch's detected-frequency set.
\end{theorem}

\begin{theorem}[Constancy after reindexing the restriction]
\lean{FABL.juntaRestriction_eq_const_of_reindexed_eq_const}
\label{FABL.juntaRestriction_eq_const_of_reindexed_eq_const}
\leanok
\uses{FABL.BooleanFunction, FABL.JuntaFixedAssignment, FABL.Sign, FABL.fixedSignCubeEquiv, FABL.juntaRestriction, FABL.reindexedJuntaRestriction}
\difficulty{3}
Let $n\in\mathbb{N}$, let $f:\{−1,1\}^n\to\{−1,1\}$ be a Boolean function, let
$P\subseteq\{1,\ldots,n\}$, fix an assignment $z$ on $P$, and let $c\in\{−1,1\}$. If the
reindexed restriction of $f$ is constantly $c$ on the sign cube whose dimension is the
number of coordinates outside $P$, then the restriction itself is constantly $c$ on
assignments to those free coordinates.
\end{theorem}

\begin{theorem}[Low degree junta node decision from batch is correct of first detected eq some]
\lean{FABL.lowDegreeJuntaNodeDecisionFromBatch_isCorrect_of_firstDetected_eq_some}
\statementsource{Analysis-of-Boolean-Functions-by-Ryan-ODonnell}{
  pdf-ffbe9c7b3306-p083-b002,
  pdf-ffbe9c7b3306-p156-b006,
  pdf-ffbe9c7b3306-p168-b001
}
\label{FABL.lowDegreeJuntaNodeDecisionFromBatch_isCorrect_of_firstDetected_eq_some}
\leanok
\uses{FABL.BooleanFunction, FABL.BooleanFunction.toReal, FABL.IsRelevant, FABL.JuntaFixedAssignment, FABL.JuntaFreeAssignment, FABL.JuntaNodeDecision.IsCorrect, FABL.LowDegreeJuntaFourierEstimatesAccurate, FABL.PositiveLearningParameter, FABL.Sign, FABL.firstLowDegreeJuntaDetectedFrequency, FABL.fourierCoeff_ne_zero_of_mem_lowDegreeJuntaDetectedFrequencies, FABL.isRelevantJuntaRestriction_of_reindexed, FABL.isRelevant_toReal_of_fourierCoeff_ne_zero, FABL.juntaRestrictionDimension, FABL.lowDegreeJuntaNodeDecisionFromBatch, FABL.lowDegreeJuntaRestrictedSampleCount, FABL.mem_detectedFrequencies_of_firstDetected_eq_some, FABL.reindexedJuntaRestriction}
\difficulty{4}
% dataset-pipeline: stub -- prose pending build_informal_statements.py
Statement pending regeneration.
\end{theorem}

\begin{theorem}[Low degree junta node decision from batch is correct]
\lean{FABL.lowDegreeJuntaNodeDecisionFromBatch_isCorrect}
\statementsource{Analysis-of-Boolean-Functions-by-Ryan-ODonnell}{
  pdf-ffbe9c7b3306-p156-b006,
  pdf-ffbe9c7b3306-p156-b007,
  pdf-ffbe9c7b3306-p168-b001
}
\label{FABL.lowDegreeJuntaNodeDecisionFromBatch_isCorrect}
\leanok
\uses{FABL.BooleanFunction, FABL.BooleanFunction.toReal, FABL.IsRelevant, FABL.JuntaFixedAssignment, FABL.JuntaFreeAssignment, FABL.JuntaNodeDecision.IsCorrect, FABL.LowDegreeF₂Hypothesis.evaluate, FABL.LowDegreeJuntaFourierEstimatesAccurate, FABL.LowDegreeJuntaPolynomialHypothesisExact, FABL.PositiveLearningParameter, FABL.Sign, FABL.binaryCubeSignEquiv, FABL.binarySignEquiv, FABL.firstLowDegreeJuntaDetectedFrequency, FABL.firstNonconstantLowDegreeMonomial, FABL.firstNonconstantLowDegreeMonomial_sound, FABL.f₂HypothesisBooleanFunction, FABL.f₂HypothesisBooleanFunction_encoding, FABL.isRelevantJuntaRestriction_of_reindexed, FABL.juntaFreeCoordinateOfFin, FABL.juntaRestriction_eq_const_of_reindexed_eq_const, FABL.lowDegreeF₂Eval, FABL.lowDegreeJuntaNodeDecisionFromBatch, FABL.lowDegreeJuntaNodeDecisionFromBatch_isCorrect_of_firstDetected_eq_some, FABL.lowDegreeJuntaPolynomialHypothesis, FABL.lowDegreeJuntaRestrictedSampleCount, FABL.reindexedJuntaRestriction}
\difficulty{5}
% dataset-pipeline: stub -- prose pending build_informal_statements.py
Statement pending regeneration.
\end{theorem}

\begin{definition}[Bad junta node batch]
\lean{FABL.LowDegreeJuntaNodeBatchBad}
\statementsource{Analysis-of-Boolean-Functions-by-Ryan-ODonnell}{pdf-ffbe9c7b3306-p168-b001}
\label{FABL.LowDegreeJuntaNodeBatchBad}
\leanok
\uses{FABL.BooleanFunction, FABL.JuntaFixedAssignment, FABL.JuntaFreeAssignment, FABL.JuntaNodeDecision.IsCorrect, FABL.PositiveLearningParameter, FABL.Sign, FABL.lowDegreeJuntaNodeDecisionFromBatch, FABL.lowDegreeJuntaRestrictedSampleCount}
Let $n$ be a natural number, let $f$ be a Boolean function on the $n$-dimensional sign
cube, let $P$ be a set of coordinates, and let $z$ assign signs to the coordinates in
$P$. Fix a natural number $k$, a failure parameter in $(0,1/2]$, and a batch of the
specified restricted-sample size, each sample consisting of an assignment on the
coordinates outside $P$ together with a sign. The batch is bad exactly when the node
decision computed from it is not correct for the restriction of $f$ obtained by fixing
$P$ according to $z$.
\end{definition}

\begin{definition}[Bad junta node source]
\lean{FABL.lowDegreeJuntaNodeSourceBad}
\statementsource{Analysis-of-Boolean-Functions-by-Ryan-ODonnell}{pdf-ffbe9c7b3306-p168-b001}
\label{FABL.lowDegreeJuntaNodeSourceBad}
\leanok
\uses{FABL.BooleanFunction, FABL.JuntaFixedAssignment, FABL.JuntaRestrictionSampleBad, FABL.LowDegreeJuntaNodeBatchBad, FABL.MatchedJuntaExample, FABL.PositiveLearningParameter, FABL.lowDegreeJuntaRestrictedSampleCount}
Let $n$ be a natural number, let $f$ be a Boolean function on the $n$-dimensional sign
cube, let $P$ be a set of coordinates, and let $z$ assign signs to the coordinates in
$P$. Fix a natural number $k$, a failure parameter in $(0,1/2]$, and either no sample
collection or a collection of the specified restricted-sample size whose entries are
cube points and signs matching $z$ on $P$. This source is bad exactly when the
restriction-sample badness predicate holds for this optional collection and the event
that the corresponding batch yields an incorrect node decision.
\end{definition}

\begin{theorem}[Ideal low degree junta node batch bad probability le]
\lean{FABL.idealLowDegreeJuntaNodeBatchBadProbability_le}
\statementsource{Analysis-of-Boolean-Functions-by-Ryan-ODonnell}{
  pdf-ffbe9c7b3306-p080-b003,
  pdf-ffbe9c7b3306-p085-b004,
  pdf-ffbe9c7b3306-p156-b006,
  pdf-ffbe9c7b3306-p156-b007,
  pdf-ffbe9c7b3306-p168-b001
}
\label{FABL.idealLowDegreeJuntaNodeBatchBadProbability_le}
\leanok
\uses{FABL.BooleanFunction, FABL.BooleanFunction.toReal, FABL.IsCorrelationImmune, FABL.JuntaFixedAssignment, FABL.JuntaFreeAssignment, FABL.JuntaNodeDecision.IsCorrect, FABL.LowDegreeJuntaFourierEstimatesAccurate, FABL.LowDegreeJuntaNodeBatchBad, FABL.LowDegreeJuntaPolynomialHypothesisExact, FABL.PositiveLearningParameter, FABL.Sign, FABL.booleanFunctionF₂Encoding, FABL.firstLowDegreeJuntaDetectedFrequency, FABL.fourierCoeff, FABL.fourierCoeff_eq_zero_of_no_lowDegreeJuntaDetection, FABL.functionAlgebraicDegree, FABL.functionAlgebraicDegree_encoding_le_sub_of_isKJunta_of_isCorrelationImmune, FABL.idealLowDegreeJuntaFourierInaccuracyProbability_le, FABL.idealLowDegreeJuntaPolynomialFailureProbability_le, FABL.juntaRestrictionLabeledBatch, FABL.lowDegreeJuntaCutoff, FABL.lowDegreeJuntaFourierFamily, FABL.lowDegreeJuntaHalfFailure, FABL.lowDegreeJuntaHalfFailure_value, FABL.lowDegreeJuntaNodeDecisionFromBatch, FABL.lowDegreeJuntaNodeDecisionFromBatch_isCorrect, FABL.lowDegreeJuntaNodeDecisionFromBatch_isCorrect_of_firstDetected_eq_some, FABL.lowDegreeJuntaPolynomialDegree, FABL.lowDegreeJuntaQuarterFailure, FABL.lowDegreeJuntaQuarterFailure_value, FABL.lowDegreeJuntaRestrictedSampleCount, FABL.mem_lowDegreeJuntaFourierFamily, FABL.reindexedJuntaRestriction, FABL.reindexedJuntaRestriction_isKJunta, FABL.uniformPMF}
\difficulty{6}
% dataset-pipeline: stub -- prose pending build_informal_statements.py
Statement pending regeneration.
\end{theorem}

\begin{theorem}[Correctness of the Junta Node Continuation]
\lean{FABL.lowDegreeJuntaNodeContinuation_badProbability_le_zero}
\label{FABL.lowDegreeJuntaNodeContinuation_badProbability_le_zero}
\leanok
\uses{FABL.BooleanFunction, FABL.JuntaFixedAssignment, FABL.JuntaNodeDecision.IsBad, FABL.JuntaNodeDecision.IsCorrect, FABL.JuntaRestrictionSampleBad, FABL.LearningProgram.addOutcomeCost, FABL.LearningProgram.eventProbability, FABL.LearningProgram.runWithCost, FABL.LowDegreeJuntaNodeBatchBad, FABL.MatchedJuntaExample, FABL.PositiveLearningParameter, FABL.lowDegreeJuntaNodeAnalysisWork, FABL.lowDegreeJuntaNodeContinuation, FABL.lowDegreeJuntaNodeDecisionFromBatch, FABL.lowDegreeJuntaNodeSourceBad, FABL.lowDegreeJuntaRestrictedSampleCount, FABL.matchedJuntaRestrictionExample, FABL.projectJuntaRestrictionBatch}
\difficulty{4}
Let $n,k\in\mathbb N$, let $f:\{\pm1\}^n\to\{\pm1\}$ be a Boolean function, let
$P\subseteq[n]$, and fix a sign assignment $z:P\to\{\pm1\}$. Let $\delta\in\mathbb Q$
satisfy $0<\delta\le 1/2$, and let the sampled input be either absent or a batch of the
prescribed size whose examples agree with $z$ on $P$. If this sampled input is not bad
for the node procedure, then the probability that the continuation returns a bad node
decision is at most zero.
\end{theorem}

\begin{theorem}[Low degree junta relevant coordinate program failure probability le]
\lean{FABL.lowDegreeJuntaRelevantCoordinateProgram_failureProbability_le}
\statementsource{Analysis-of-Boolean-Functions-by-Ryan-ODonnell}{
  pdf-ffbe9c7b3306-p156-b005,
  pdf-ffbe9c7b3306-p168-b001
}
\label{FABL.lowDegreeJuntaRelevantCoordinateProgram_failureProbability_le}
\leanok
\uses{FABL.BooleanFunction, FABL.JuntaFixedAssignment, FABL.JuntaNodeDecision, FABL.JuntaNodeDecision.IsBad, FABL.JuntaRestrictionSampleBad, FABL.LearningProgram, FABL.LearningProgram.bind, FABL.LearningProgram.eventProbability, FABL.LearningProgram.eventProbability_bind_le_add, FABL.LowDegreeJuntaNodeBatchBad, FABL.MatchedJuntaExample, FABL.PositiveLearningParameter, FABL.idealLowDegreeJuntaNodeBatchBadProbability_le, FABL.juntaRestrictionSampleProgram, FABL.juntaRestrictionSampleProgram_badProbability_le, FABL.lowDegreeJuntaHalfFailure, FABL.lowDegreeJuntaHalfFailure_value, FABL.lowDegreeJuntaNodeContinuation, FABL.lowDegreeJuntaNodeContinuation_badProbability_le_zero, FABL.lowDegreeJuntaNodeSourceBad, FABL.lowDegreeJuntaRelevantCoordinateProgram, FABL.lowDegreeJuntaRestrictedSampleCount, FABL.positiveLearningParameter_toReal_mem_Ioc}
\difficulty{5}
% dataset-pipeline: stub -- prose pending build_informal_statements.py
Statement pending regeneration.
\end{theorem}

\begin{definition}[Bad root-finder output]
\lean{FABL.LowDegreeJuntaRootFinderBad}
\statementsource{Analysis-of-Boolean-Functions-by-Ryan-ODonnell}{pdf-ffbe9c7b3306-p168-b001}
\label{FABL.LowDegreeJuntaRootFinderBad}
\leanok
\uses{FABL.BooleanFunction, FABL.BooleanFunction.toReal, FABL.IsRelevant, FABL.JuntaNodeDecision}
Let $n$ be a natural number and let $f$ be a Boolean function on the $n$-dimensional
sign cube. An optional root decision is declared bad as follows: absence of a decision
is not bad, and a constant-sign decision is not bad; a decision to split on a coordinate
is bad exactly when that coordinate has zero influence on $f$, where influence is the
mean squared discrete derivative under the uniform sign-cube distribution.
\end{definition}

\begin{theorem}[Restriction with no fixed coordinates]
\lean{FABL.juntaRestriction_empty_juntaFreePart}
\label{FABL.juntaRestriction_empty_juntaFreePart}
\leanok
\uses{FABL.BooleanFunction, FABL.JuntaFixedAssignment, FABL.MatchesJuntaAssignment, FABL.combineJuntaAssignment_freePart_of_matches, FABL.juntaFreePart, FABL.juntaRestriction}
\difficulty{3}
Let $n\in\mathbb{N}$, let $f:\{−1,1\}^n\to\{−1,1\}$ be a Boolean function, let $z$ be an
assignment on the empty coordinate set, and let $x\in\{−1,1\}^n$. The restriction of $f$
with no coordinates fixed, evaluated at the free-coordinate part of $x$, equals $f(x)$.
\end{theorem}

\begin{theorem}[Relevance of the Unrestricted Function]
\lean{FABL.isRelevant_of_isRelevantJuntaRestriction_empty}
\label{FABL.isRelevant_of_isRelevantJuntaRestriction_empty}
\leanok
\uses{FABL.BooleanFunction, FABL.BooleanFunction.toReal, FABL.IsRelevant, FABL.IsRelevantJuntaRestriction, FABL.JuntaFixedAssignment, FABL.JuntaFreeIndex, FABL.Sign, FABL.combineJuntaAssignment, FABL.combineJuntaAssignment_apply_free, FABL.isRelevant_iff_exists_setCoordinate_ne, FABL.setCoordinate, FABL.setCoordinate_apply_of_ne, FABL.setCoordinate_apply_self, FABL.setCoordinate_eq_self, FABL.signValue_injective}
\difficulty{6}
Let $f$ be a Boolean function on the $n$-dimensional sign cube, and let $i$ be a
coordinate. If there are two sign-cube assignments that differ only at $i$ and on which
$f$ takes different values, then $f$, viewed as a real-valued function, has positive
influence at $i$.
\end{theorem}

\begin{theorem}[Root-finder failure gives a bad empty-set decision]
\lean{FABL.juntaNodeDecision_isBad_empty_of_rootFinderBad}
\label{FABL.juntaNodeDecision_isBad_empty_of_rootFinderBad}
\leanok
\uses{FABL.BooleanFunction, FABL.BooleanFunction.toReal, FABL.IsRelevant, FABL.IsRelevantJuntaRestriction, FABL.JuntaNodeDecision, FABL.JuntaNodeDecision.IsBad, FABL.LowDegreeJuntaRootFinderBad, FABL.isRelevant_of_isRelevantJuntaRestriction_empty, FABL.juntaFreePart, FABL.juntaRestriction, FABL.juntaRestriction_empty_juntaFreePart}
\difficulty{4}
Let $n$ be a natural number and let $f$ be a Boolean function on sign vectors of length
$n$ that is not constant. Let an outcome be either absent or a decision for the
restriction with no fixed coordinates. If the root-finder outcome is bad—meaning it is
absent, is a constant decision, or, when it selects a coordinate, that coordinate has
zero influence on $f$—then the outcome is bad as a decision for the empty-set
restriction with the fixed assignment taking value $1$ wherever required.
\end{theorem}

\begin{theorem}[Low degree junta relevant coordinate program root failure probability le]
\lean{FABL.lowDegreeJuntaRelevantCoordinateProgram_root_failureProbability_le}
\statementsource{Analysis-of-Boolean-Functions-by-Ryan-ODonnell}{
  pdf-ffbe9c7b3306-p156-b005,
  pdf-ffbe9c7b3306-p168-b001
}
\label{FABL.lowDegreeJuntaRelevantCoordinateProgram_root_failureProbability_le}
\leanok
\uses{FABL.BooleanFunction, FABL.IsKJunta, FABL.LearningProgram.eventProbability, FABL.LearningProgram.eventProbability_mono, FABL.LowDegreeJuntaRootFinderBad, FABL.PositiveLearningParameter, FABL.juntaNodeDecision_isBad_empty_of_rootFinderBad, FABL.lowDegreeJuntaRelevantCoordinateProgram, FABL.lowDegreeJuntaRelevantCoordinateProgram_failureProbability_le}
\difficulty{3}
% dataset-pipeline: stub -- prose pending build_informal_statements.py
Statement pending regeneration.
\end{theorem}

\begin{theorem}[Low degree junta relevant coordinate program cost le]
\lean{FABL.lowDegreeJuntaRelevantCoordinateProgram_cost_le}
\statementsource{Analysis-of-Boolean-Functions-by-Ryan-ODonnell}{
  pdf-ffbe9c7b3306-p156-b005,
  pdf-ffbe9c7b3306-p156-b006,
  pdf-ffbe9c7b3306-p156-b007,
  pdf-ffbe9c7b3306-p156-b008
}
\label{FABL.lowDegreeJuntaRelevantCoordinateProgram_cost_le}
\leanok
\uses{FABL.BooleanFunction, FABL.JuntaFixedAssignment, FABL.JuntaNodeDecision, FABL.LearningCost, FABL.LearningProgram, FABL.LearningProgram.bind, FABL.LearningProgram.runWithCost, FABL.LearningProgram.runWithCost_bind, FABL.MatchedJuntaExample, FABL.PositiveLearningParameter, FABL.juntaRestrictionSampleCount, FABL.juntaRestrictionSampleProgram, FABL.juntaRestrictionSampleProgram_cost_eq, FABL.juntaRestrictionSampleWork, FABL.lowDegreeJuntaHalfFailure, FABL.lowDegreeJuntaNodeAnalysisWork, FABL.lowDegreeJuntaNodeAnalysisWorkBound, FABL.lowDegreeJuntaNodeAnalysisWork_le, FABL.lowDegreeJuntaNodeContinuation, FABL.lowDegreeJuntaNodeRandomExampleBound, FABL.lowDegreeJuntaNodeRandomExamples, FABL.lowDegreeJuntaNodeUniformWorkBound, FABL.lowDegreeJuntaNodeWorkBound, FABL.lowDegreeJuntaRelevantCoordinateProgram, FABL.lowDegreeJuntaRestrictedSampleCount, FABL.matchedJuntaRestrictionExample}
\difficulty{5}
% dataset-pipeline: stub -- prose pending build_informal_statements.py
Statement pending regeneration.
\end{theorem}

\begin{definition}[Low-degree relevant-coordinate finder]
\lean{FABL.lowDegreeJuntaRelevantCoordinateFinder}
\statementsource{Analysis-of-Boolean-Functions-by-Ryan-ODonnell}{
  pdf-ffbe9c7b3306-p156-b005,
  pdf-ffbe9c7b3306-p156-b008,
  pdf-ffbe9c7b3306-p168-b001
}
\label{FABL.lowDegreeJuntaRelevantCoordinateFinder}
\leanok
\uses{FABL.JuntaRelevantCoordinateFinder, FABL.lowDegreeJuntaNodeRandomExampleBound, FABL.lowDegreeJuntaNodeUniformWorkBound, FABL.lowDegreeJuntaRelevantCoordinateProgram, FABL.lowDegreeJuntaRelevantCoordinateProgram_cost_le, FABL.lowDegreeJuntaRelevantCoordinateProgram_failureProbability_le}
For natural numbers $n$ and $k$, a low-degree relevant-coordinate finder assigns to each
coordinate set $P\subseteq\{0,\ldots,n-1\}$, each fixed sign assignment $z$ on $P$, and
each failure parameter in $(0,1/2]$ a randomized program using random labeled examples
and returning either no decision or a constant-sign or coordinate-split decision. It
also specifies natural-number bounds on the number of random-example calls and on local
work. Every outcome in the program's support uses at most the declared number of
examples, makes zero membership queries, and uses at most the declared local work;
moreover, for every target depending only on some coordinate set $J$ of size at most
$k$, whenever $P\subseteq J$, the probability that the returned decision is absent or
incorrect for the restriction fixed by $z$ is at most the given failure parameter.
\end{definition}

\begin{definition}[Recursive junta learner program]
\lean{FABL.lowDegreeJuntaLearnerProgram}
\statementsource{Analysis-of-Boolean-Functions-by-Ryan-ODonnell}{
  pdf-ffbe9c7b3306-p081-b003,
  pdf-ffbe9c7b3306-p156-b002,
  pdf-ffbe9c7b3306-p156-b005,
  pdf-ffbe9c7b3306-p168-b001
}
\label{FABL.lowDegreeJuntaLearnerProgram}
\leanok
\uses{FABL.DecisionTree, FABL.LearningProgram, FABL.PositiveLearningParameter, FABL.Sign, FABL.lowDegreeJuntaRelevantCoordinateFinder, FABL.recursiveJuntaLearner}
For natural numbers $n$ and $k$ and a failure parameter in $(0,1/2]$, the learner
program is the recursive junta learner built using the low-degree relevant-coordinate
finder for parameters $n$ and $k$. It uses random-example access and returns either no
decision tree or a decision tree on the $n$-dimensional sign cube with sign-valued
leaves.
\end{definition}

\begin{theorem}[Low degree junta learner program failure probability le]
\lean{FABL.lowDegreeJuntaLearnerProgram_failureProbability_le}
\statementsource{Analysis-of-Boolean-Functions-by-Ryan-ODonnell}{
  pdf-ffbe9c7b3306-p156-b005,
  pdf-ffbe9c7b3306-p168-b001,
  pdf-ffbe9c7b3306-p081-b003
}
\label{FABL.lowDegreeJuntaLearnerProgram_failureProbability_le}
\leanok
\uses{FABL.BooleanFunction, FABL.IsKJunta, FABL.JuntaLearnerOutputBad, FABL.LearningProgram.eventProbability, FABL.PositiveLearningParameter, FABL.lowDegreeJuntaLearnerProgram, FABL.lowDegreeJuntaRelevantCoordinateFinder, FABL.recursiveJuntaLearner_failureProbability_le}
\difficulty{1}
% dataset-pipeline: stub -- prose pending build_informal_statements.py
Statement pending regeneration.
\end{theorem}

\begin{theorem}[Depth of the junta learner's output]
\lean{FABL.lowDegreeJuntaLearnerProgram_depth_le}
\proofsource{Analysis-of-Boolean-Functions-by-Ryan-ODonnell}{pdf-ffbe9c7b3306-p168-b001}
\label{FABL.lowDegreeJuntaLearnerProgram_depth_le}
\leanok
\uses{FABL.BooleanFunction, FABL.DecisionTree, FABL.F₂DecisionTree.depth, FABL.LearningCost, FABL.LearningProgram.runWithCost, FABL.PositiveLearningParameter, FABL.Sign, FABL.lowDegreeJuntaLearnerProgram, FABL.lowDegreeJuntaRelevantCoordinateFinder, FABL.recursiveJuntaLearner_depth_le}
\difficulty{1}
Let $n,k\in\mathbb{N}$, let $f:\{\pm1\}^n\to\{\pm1\}$ be a Boolean function, and let
$\delta\in\mathbb{Q}$ satisfy $0<\delta\le 1/2$. If an outcome lies in the support of
the learner's run on $f$ and its output is the decision tree $T$, then $T$ has depth at
most $k$.
\end{theorem}

\begin{theorem}[Learner program cost bound]
\lean{FABL.lowDegreeJuntaLearnerProgram_cost_le}
\proofsource{Analysis-of-Boolean-Functions-by-Ryan-ODonnell}{
  pdf-ffbe9c7b3306-p156-b002,
  pdf-ffbe9c7b3306-p168-b001
}
\label{FABL.lowDegreeJuntaLearnerProgram_cost_le}
\leanok
\uses{FABL.BooleanFunction, FABL.DecisionTree, FABL.LearningCost, FABL.LearningProgram.runWithCost, FABL.PositiveLearningParameter, FABL.Sign, FABL.juntaTreeCallCount, FABL.juntaTreePerCallFailure, FABL.lowDegreeJuntaLearnerProgram, FABL.lowDegreeJuntaNodeRandomExampleBound, FABL.lowDegreeJuntaNodeUniformWorkBound, FABL.lowDegreeJuntaRelevantCoordinateFinder, FABL.recursiveJuntaLearner_cost_le}
\difficulty{1}
Let $n,k$ be natural numbers, let $f:\{−1,1\}^n\to\{−1,1\}$ be a Boolean function, and
let $\delta\in\mathbb{Q}$ satisfy $0<\delta\le 1/2$. For every outcome in the support of
the learner's run on $f$, the number of random examples is at most $T_k
R(n,k,\delta/T_k)$, the number of membership queries is zero, and the local work is at
most $T_k W(n,k,\delta/T_k)$, where $T_0=1$, $T_{j+1}=1+2T_j$, and $R(n,k,q)$ and
$W(n,k,q)$ are respectively the learner's uniform per-node random-example and local-work
bounds at parameters $n,k,q$.
\end{theorem}

\begin{theorem}[Dimension of the free coordinates]
\lean{FABL.juntaRestrictionDimension_le}
\label{FABL.juntaRestrictionDimension_le}
\leanok
\uses{FABL.juntaRestrictionDimension}
\difficulty{3}
For every natural number $n$ and every finite set $P$ of coordinates among $n$, the
number of coordinates outside $P$ is at most $n$.
\end{theorem}

\begin{theorem}[Bound for a restricted Fourier family]
\lean{FABL.card_lowDegreeJuntaFourierFamily_le}
\proofsource{Analysis-of-Boolean-Functions-by-Ryan-ODonnell}{pdf-ffbe9c7b3306-p081-b004}
\label{FABL.card_lowDegreeJuntaFourierFamily_le}
\leanok
\uses{FABL.card_lowDegreeFourierFamily_le, FABL.juntaRestrictionDimension, FABL.juntaRestrictionDimension_le, FABL.lowDegreeJuntaCutoff, FABL.lowDegreeJuntaFourierFamily}
\difficulty{4}
Let $n$ and $k$ be natural numbers, and let $P$ be a subset of an $n$-element set. Write
$d$ for the number of elements outside $P$, and set $r=\lfloor 3k/4\rfloor$. The number
of subsets of a $d$-element set having size at most $r$ is at most $(r+1)(n+1)^r$.
\end{theorem}

\begin{theorem}[Coefficient confidence failure-bit identity]
\lean{FABL.lowDegreeJuntaCoefficientFailureBits_eq}
\label{FABL.lowDegreeJuntaCoefficientFailureBits_eq}
\leanok
\uses{FABL.PositiveLearningParameter, FABL.fourierEstimatorFailureBits, FABL.lowDegreeJuntaCoefficientConfidence, FABL.lowDegreeJuntaCoefficientConfidence_value, FABL.lowDegreeJuntaFourierFamily, FABL.lowDegreeJuntaFourierFamily_nonempty}
\difficulty{4}
Let $n$ be a natural number, let $P\subseteq[n]$, let $k$ be a natural, and let
$\delta\in\mathbb{Q}$ satisfy $0<\delta\leq1/2$. Let $\mathcal F$ be the family of
subsets of the coordinates outside $P$ having size at most $\lfloor3k/4\rfloor$. The
number of failure bits at confidence $\delta/(4|\mathcal F|)$ is exactly
$\lceil\log_2\!\left(\left\lceil 8|\mathcal F|/\delta\right\rceil\right)\rceil$.
\end{theorem}

\begin{theorem}[Fourier Scan Work Bound]
\lean{FABL.lowDegreeJuntaFourierScanWork_cast_le}
\label{FABL.lowDegreeJuntaFourierScanWork_cast_le}
\leanok
\uses{FABL.PositiveLearningParameter, FABL.card_lowDegreeJuntaFourierFamily_le, FABL.fourierEstimatorFailureBits, FABL.fourierEstimatorSampleCount_cast_le, FABL.lowDegreeJuntaCoefficientAccuracy, FABL.lowDegreeJuntaCoefficientConfidence, FABL.lowDegreeJuntaCoefficientSampleCount, FABL.lowDegreeJuntaCutoff, FABL.lowDegreeJuntaFourierFamily}
\difficulty{4}
Let $n,k$ be natural numbers, let $P$ be a subset of an $n$-element index set, and let
$\delta$ be rational with $0<\delta\le1/2$. Put $d=n-|P|$, $r=\lfloor3k/4\rfloor$, and
let $\mathcal F$ be the subsets of a $d$-element set of size at most $r$. Define
$\alpha=1/(3\cdot2^k)$ and $\eta=\delta/(4|\mathcal F|T_k)$, where $T_0=1$ and
$T_{j+1}=1+2T_j$. For $q>0$, set $B(q)=\lceil\log_2\lceil2/q\rceil\rceil$ and
$S=\lceil2B(\eta)/\alpha^2\rceil$. Then $|\mathcal F|Sr\le
(r+1)(n+1)^r\bigl(4B(\eta)/\alpha^2\bigr)(r+1)$.
\end{theorem}

\begin{theorem}[Low degree junta polynomial work le fixed parameter envelope]
\lean{FABL.lowDegreeJuntaPolynomialWork_le_fixedParameterEnvelope}
\statementsource{Analysis-of-Boolean-Functions-by-Ryan-ODonnell}{pdf-ffbe9c7b3306-p167-b005}
\label{FABL.lowDegreeJuntaPolynomialWork_le_fixedParameterEnvelope}
\leanok
\uses{FABL.JuntaFreeAssignment, FABL.PositiveLearningParameter, FABL.Sign, FABL.fourierEstimatorFailureBits, FABL.juntaRestrictionDimension, FABL.juntaRestrictionDimension_le, FABL.lowDegreeF₂PolynomialLearnerWork, FABL.lowDegreeJuntaPolynomialDegree, FABL.lowDegreeJuntaPolynomialSampleCount, FABL.lowDegreeJuntaPolynomialSamples, FABL.lowDegreeJuntaQuarterFailure, FABL.lowDegreeJuntaRestrictedSampleCount, FABL.scheduledLowDegreeF₂PolynomialLearnerWork_le_fixedParameterEnvelope}
\difficulty{4}
% dataset-pipeline: stub -- prose pending build_informal_statements.py
Statement pending regeneration.
\end{theorem}

\begin{theorem}[Junta Runtime Exponent Inequalities]
\lean{FABL.lowDegreeJuntaRuntimeExponentLedger}
\label{FABL.lowDegreeJuntaRuntimeExponentLedger}
\leanok
\uses{FABL.lowDegreeJuntaCutoff, FABL.lowDegreeJuntaPolynomialDegree, FABL.three_mul_lowDegreeJuntaPolynomialDegree_le_cutoff_add_three}
\difficulty{1}
For every $k\in\mathbb{N}$, put $q=\lfloor3k/4\rfloor$ and $\ell=k-q$. Then $q\le q+3$
and $3\ell\le q+3$.
\end{theorem}

\begin{theorem}[Low degree junta node runtime ledger]
\lean{FABL.lowDegreeJuntaNodeRuntimeLedger}
\statementsource{Analysis-of-Boolean-Functions-by-Ryan-ODonnell}{
  pdf-ffbe9c7b3306-p156-b006,
  pdf-ffbe9c7b3306-p156-b007,
  pdf-ffbe9c7b3306-p156-b008
}
\label{FABL.lowDegreeJuntaNodeRuntimeLedger}
\leanok
\uses{FABL.JuntaFreeAssignment, FABL.PositiveLearningParameter, FABL.Sign, FABL.fourierEstimatorFailureBits, FABL.juntaRestrictionDimension, FABL.lowDegreeF₂PolynomialLearnerWork, FABL.lowDegreeJuntaCoefficientAccuracy, FABL.lowDegreeJuntaCoefficientConfidence, FABL.lowDegreeJuntaCoefficientFailureBits_eq, FABL.lowDegreeJuntaCoefficientSampleCount, FABL.lowDegreeJuntaCutoff, FABL.lowDegreeJuntaFourierFamily, FABL.lowDegreeJuntaFourierScanWork_cast_le, FABL.lowDegreeJuntaPolynomialDegree, FABL.lowDegreeJuntaPolynomialSampleCount, FABL.lowDegreeJuntaPolynomialSamples, FABL.lowDegreeJuntaPolynomialWork_le_fixedParameterEnvelope, FABL.lowDegreeJuntaQuarterFailure, FABL.lowDegreeJuntaRestrictedSampleCount, FABL.lowDegreeJuntaRuntimeExponentLedger}
\difficulty{1}
% dataset-pipeline: stub -- prose pending build_informal_statements.py
Statement pending regeneration.
\end{theorem}

\begin{definition}[Positive rational exponents]
\lean{FABL.PositiveMatrixExponent}
\label{FABL.PositiveMatrixExponent}
\leanok
The positive matrix exponents are the positive rational numbers.
\end{definition}

\begin{theorem}[Positivity after adding one]
\lean{FABL.positiveMatrixExponent_add_one_pos}
\label{FABL.positiveMatrixExponent_add_one_pos}
\leanok
\uses{FABL.PositiveMatrixExponent}
\difficulty{3}
If $\\omega$ is a positive rational number, then $\\omega+1>0$.
\end{theorem}

\begin{definition}[Balanced junta exponent]
\lean{FABL.matrixExponentJuntaBalancedExponent}
\statementsource{Analysis-of-Boolean-Functions-by-Ryan-ODonnell}{
  pdf-ffbe9c7b3306-p156-b003,
  pdf-ffbe9c7b3306-p156-b008,
  pdf-ffbe9c7b3306-p167-b005
}
\label{FABL.matrixExponentJuntaBalancedExponent}
\leanok
\uses{FABL.PositiveMatrixExponent}
For a positive rational number $\omega$ and a natural number $k$, the balanced exponent
is $\omega k/(\omega+1)$.
\end{definition}

\begin{definition}[Junta exponent cutoff]
\lean{FABL.matrixExponentJuntaCutoff}
\statementsource{Analysis-of-Boolean-Functions-by-Ryan-ODonnell}{
  pdf-ffbe9c7b3306-p156-b008,
  pdf-ffbe9c7b3306-p167-b005
}
\label{FABL.matrixExponentJuntaCutoff}
\leanok
\uses{FABL.PositiveMatrixExponent, FABL.matrixExponentJuntaBalancedExponent}
For a positive rational number $\omega$ and a natural number $k$, the cutoff is the
ceiling of the rational number $\omega k/(\omega+1)$.
\end{definition}

\begin{definition}[Matrix exponent junta solver exponent]
\lean{FABL.matrixExponentJuntaSolverExponent}
\statementsource{Analysis-of-Boolean-Functions-by-Ryan-ODonnell}{
  pdf-ffbe9c7b3306-p156-b003,
  pdf-ffbe9c7b3306-p156-b008,
  pdf-ffbe9c7b3306-p167-b005
}
\label{FABL.matrixExponentJuntaSolverExponent}
\leanok
\uses{FABL.PositiveMatrixExponent, FABL.matrixExponentJuntaCutoff}
For a positive rational exponent $\omega$ and a natural number $k$, define the matrix
exponent junta solver exponent to be $\omega$ multiplied by $k$ minus the matrix
exponent junta cutoff at $\omega$ and $k$; explicitly,
$\omega\bigl(k-\operatorname{cutoff}(\omega,k)\bigr)$.
\end{definition}

\begin{theorem}[Nonnegativity of the balanced exponent]
\lean{FABL.matrixExponentJuntaBalancedExponent_nonneg}
\label{FABL.matrixExponentJuntaBalancedExponent_nonneg}
\leanok
\uses{FABL.PositiveMatrixExponent, FABL.matrixExponentJuntaBalancedExponent, FABL.positiveMatrixExponent_add_one_pos}
\difficulty{3}
Let $[?]$ be a positive rational matrix exponent and let $k$ be a natural number. The
balanced exponent is nonnegative.
\end{theorem}

\begin{theorem}[Balanced exponent bounded by degree]
\lean{FABL.matrixExponentJuntaBalancedExponent_le_k}
\label{FABL.matrixExponentJuntaBalancedExponent_le_k}
\leanok
\uses{FABL.PositiveMatrixExponent, FABL.matrixExponentJuntaBalancedExponent, FABL.positiveMatrixExponent_add_one_pos}
\difficulty{3}
Let $[?]$ be a positive rational matrix exponent and let $k$ be a natural number. The
balanced exponent is at most $k$.
\end{theorem}

\begin{theorem}[Balanced exponent identity]
\lean{FABL.matrixExponentJuntaBalancedExponent_balance}
\proofsource{Analysis-of-Boolean-Functions-by-Ryan-ODonnell}{pdf-ffbe9c7b3306-p156-b008}
\label{FABL.matrixExponentJuntaBalancedExponent_balance}
\leanok
\uses{FABL.PositiveMatrixExponent, FABL.matrixExponentJuntaBalancedExponent, FABL.positiveMatrixExponent_add_one_pos}
\difficulty{3}
Let $[?]$ be a positive rational matrix exponent and let $k$ be a natural number. The
balanced exponent $[?]k/([?]+1)$ satisfies $[?](k-[?]k/([?]+1))=[?]k/([?]+1)$.
\end{theorem}

\begin{theorem}[Balanced exponent as a fraction]
\lean{FABL.matrixExponentJuntaBalancedExponent_eq_fraction}
\label{FABL.matrixExponentJuntaBalancedExponent_eq_fraction}
\leanok
\uses{FABL.PositiveMatrixExponent, FABL.matrixExponentJuntaBalancedExponent}
\difficulty{3}
Let $[?]$ be a positive rational matrix exponent and let $k$ be a natural number. The
balanced exponent equals $([?]/([?]+1))k$.
\end{theorem}

\begin{theorem}[Balanced exponent bounded by cutoff]
\lean{FABL.matrixExponentJuntaBalancedExponent_le_cutoff}
\label{FABL.matrixExponentJuntaBalancedExponent_le_cutoff}
\leanok
\uses{FABL.PositiveMatrixExponent, FABL.matrixExponentJuntaBalancedExponent, FABL.matrixExponentJuntaCutoff}
\difficulty{1}
Let $[?]$ be a positive rational matrix exponent and let $k$ be a natural number. The
balanced exponent is at most its ceiling, the integer cutoff.
\end{theorem}

\begin{theorem}[Cutoff bounded by degree]
\lean{FABL.matrixExponentJuntaCutoff_le_k}
\label{FABL.matrixExponentJuntaCutoff_le_k}
\leanok
\uses{FABL.PositiveMatrixExponent, FABL.matrixExponentJuntaBalancedExponent_le_k, FABL.matrixExponentJuntaCutoff}
\difficulty{1}
Let $\omega$ be a positive rational number and let $k$ be a natural number. Then $\lceil
\omega k/(\omega+1)\rceil\le k$.
\end{theorem}

\begin{theorem}[Nonnegativity of the solver exponent]
\lean{FABL.matrixExponentJuntaSolverExponent_nonneg}
\label{FABL.matrixExponentJuntaSolverExponent_nonneg}
\leanok
\uses{FABL.PositiveMatrixExponent, FABL.matrixExponentJuntaCutoff, FABL.matrixExponentJuntaCutoff_le_k, FABL.matrixExponentJuntaSolverExponent}
\difficulty{3}
Let $\omega$ be a positive rational number and let $k$ be a natural number, and set
$c=\lceil \omega k/(\omega+1)\rceil$. Then $0\le\omega(k-c)$.
\end{theorem}

\begin{theorem}[Ceiling bound by balanced exponent]
\lean{FABL.matrixExponentJuntaCutoff_lt_balancedExponent_add_one}
\label{FABL.matrixExponentJuntaCutoff_lt_balancedExponent_add_one}
\leanok
\uses{FABL.PositiveMatrixExponent, FABL.matrixExponentJuntaBalancedExponent, FABL.matrixExponentJuntaBalancedExponent_nonneg, FABL.matrixExponentJuntaCutoff}
\difficulty{1}
Let $\omega$ be a positive rational number and let $k$ be a natural number. Then $\lceil
\omega k/(\omega+1)\rceil<\omega k/(\omega+1)+1$.
\end{theorem}

\begin{theorem}[Solver exponent bounded by cutoff]
\lean{FABL.matrixExponentJuntaSolverExponent_le_cutoff}
\label{FABL.matrixExponentJuntaSolverExponent_le_cutoff}
\leanok
\uses{FABL.PositiveMatrixExponent, FABL.matrixExponentJuntaBalancedExponent, FABL.matrixExponentJuntaBalancedExponent_balance, FABL.matrixExponentJuntaBalancedExponent_le_cutoff, FABL.matrixExponentJuntaCutoff, FABL.matrixExponentJuntaSolverExponent}
\difficulty{4}
Let $\omega$ be a positive rational number and let $k$ be a natural number, and set
$c=\lceil \omega k/(\omega+1)\rceil$. Then $\omega(k-c)\le c$.
\end{theorem}

\begin{theorem}[Combined exponent bound]
\lean{FABL.matrixExponentJuntaCombinedExponent_lt_fraction_add_one}
\label{FABL.matrixExponentJuntaCombinedExponent_lt_fraction_add_one}
\leanok
\uses{FABL.PositiveMatrixExponent, FABL.matrixExponentJuntaBalancedExponent_eq_fraction, FABL.matrixExponentJuntaCutoff, FABL.matrixExponentJuntaCutoff_lt_balancedExponent_add_one, FABL.matrixExponentJuntaSolverExponent, FABL.matrixExponentJuntaSolverExponent_le_cutoff}
\difficulty{3}
Let $\omega$ be a positive rational number and let $k$ be a natural number, and set
$c=\lceil \omega k/(\omega+1)\rceil$. Then
$\max\{\omega(k-c),c\}<\frac{\omega}{\omega+1}k+1$.
\end{theorem}

\begin{definition}[Cubic matrix exponent]
\lean{FABL.cubicMatrixExponent}
\label{FABL.cubicMatrixExponent}
\leanok
\uses{FABL.PositiveMatrixExponent}
The cubic matrix exponent is the positive rational number $3$.
\end{definition}

\begin{theorem}[Balanced exponent for exponent three]
\lean{FABL.matrixExponentJuntaBalancedExponent_cubic}
\proofsource{Analysis-of-Boolean-Functions-by-Ryan-ODonnell}{pdf-ffbe9c7b3306-p156-b003}
\label{FABL.matrixExponentJuntaBalancedExponent_cubic}
\leanok
\uses{FABL.cubicMatrixExponent, FABL.matrixExponentJuntaBalancedExponent}
\difficulty{3}
For every natural number $k$, the balanced exponent associated with the rational matrix
exponent $3$ equals $3k/4$.
\end{theorem}

\begin{theorem}[Matrix exponent junta runtime exponent ledger]
\lean{FABL.matrixExponentJuntaRuntimeExponentLedger}
\statementsource{Analysis-of-Boolean-Functions-by-Ryan-ODonnell}{
  pdf-ffbe9c7b3306-p156-b006,
  pdf-ffbe9c7b3306-p156-b007,
  pdf-ffbe9c7b3306-p156-b008
}
\label{FABL.matrixExponentJuntaRuntimeExponentLedger}
\leanok
\uses{FABL.PositiveMatrixExponent, FABL.matrixExponentJuntaBalancedExponent, FABL.matrixExponentJuntaBalancedExponent_eq_fraction, FABL.matrixExponentJuntaCombinedExponent_lt_fraction_add_one, FABL.matrixExponentJuntaCutoff, FABL.matrixExponentJuntaSolverExponent, FABL.matrixExponentJuntaSolverExponent_le_cutoff, FABL.matrixExponentJuntaSolverExponent_nonneg}
\difficulty{1}
% dataset-pipeline: stub -- prose pending build_informal_statements.py
Statement pending regeneration.
\end{theorem}
